\documentclass[aps,prd,twocolumn,superscriptaddress,nofootinbib]{revtex4-2}

% MiKTeX REVTeX 4.2 needs explicit package loading
\usepackage{amsmath}
\usepackage{amssymb}
\usepackage{booktabs}
\usepackage{tabularx}
\usepackage{xcolor}
\usepackage{amsthm}
\usepackage{array}
\usepackage{float}
\usepackage{graphicx}
\graphicspath{{../figures/}}

% Theorem environments (REVTeX compatible)
\newtheorem{definition}{Definition}
\newtheorem{proposition}{Proposition}
\newtheorem{conjecture}{Conjecture}

\usepackage{hyperref}
\hypersetup{
    pdftitle={Eliminating the Dark Sector: Unifying the Curvature Relaxation Model with MOND --- Replacing the Particle Dark Sector with Spacetime Geometry},
    pdfauthor={L. Geiger},
    colorlinks=true,
    linkcolor=black,
    urlcolor=blue,
    citecolor=black
}

\begin{document}

% ===================================================================
% TITLE PAGE
% ===================================================================

\title{Eliminating the Dark Sector: Unifying the Curvature Relaxation Model with MOND --- Replacing the Particle Dark Sector with Spacetime Geometry}

\author{L. Geiger}
\email{https://orcid.org/0009-0005-7296-1534}
\affiliation{Independent Researcher, Bernau, Germany}

\date{\today}

\begin{abstract}
We propose a unified geometric framework that replaces both dark energy and particulate dark matter with spacetime curvature, extending the Curvature Relaxation Model (CRM) \cite{Geiger2026} to a universe with exclusively baryonic matter ($\Omega_b \approx 0.05$), where gravitational effects attributed to dark matter arise from the scalaron of $R + 2\gamma R^2$ gravity, compatible with MOND \cite{Milgrom1983}. Three key innovations are introduced: (i)~a \textit{running curvature coupling} $\beta_{\mathrm{eff}}(a)$ transitioning from CDM-like ($\beta \approx 2.8$) to curvature-like ($\beta \approx 2.0$); (ii)~a scale-dependent MOND background coupling $\mu(a) = \sqrt{\pi}$ at late times, $\mu \to 1$ at $z > 4000$, which eliminates the need for Early Dark Energy; and (iii)~a geometric power-law term $\alpha \cdot a^{-\beta}$ replacing the cosmological role of dark matter. Tested jointly against 1,590 Pantheon+ supernovae, Planck CMB observables, and 9 BAO measurements, the preferred variant achieves $\Delta\chi^2 = -5.5$ vs.\ $\Lambda$CDM at $H_0 = 67.3$\,km/s/Mpc with zero EDE, matching Planck exactly ($\ell_A = 301.471$, $\mathcal{R} = 1.7502$), using only 6 free parameters. The SN-only analysis yields $\Delta\chi^2 = -26.3$ with MCMC posteriors $\beta = 2.02 \pm 0.20$, confirming curvature scaling. The scalaron produces testable signatures: $\mu(k,a) \neq 1$, gravitational slip $\eta \neq 1$, yet lensing $\Sigma = 1$ -- distinguishing this framework from CDM.
\end{abstract}

\keywords{Curvature Relaxation Model, MOND, dark matter, dark energy, baryon-only universe, Pantheon+, modified gravity, geometric cosmology}

\thanks{AI tools (Claude, Anthropic; Copilot, Microsoft/GitHub; Gemini, Google DeepMind) were used for mathematical formalization, code development, and text generation. All physical hypotheses, scientific interpretation, and responsibility for the content lie solely with the author. The analysis code is available at \url{https://github.com/lukisch/crm-cosmology}.}

\thanks{Paper~II in the CRM series -- Companion to \cite{Geiger2026}.}

\maketitle


% ===================================================================
% 1. INTRODUCTION
% ===================================================================
\section{Introduction: The Dark Sector Problem}
\label{sec:intro}

The standard cosmological model, $\Lambda$CDM, describes the energy budget of the universe as consisting of approximately 5\% baryonic matter, 27\% cold dark matter (CDM), and 68\% dark energy ($\Lambda$) \cite{Planck2020}. Despite its remarkable empirical success, this model implies that \textit{95\% of the universe consists of entities that have never been directly detected}.

Two independent lines of research challenge this picture:

\begin{enumerate}
\item \textbf{The Curvature Relaxation Model (CRM)} \cite{Geiger2026}: Developed from a game-theoretic framework, the CRM replaces the cosmological constant $\Lambda$ with a time-dependent curvature return potential $\Omega_\Phi(a)$, explaining accelerated expansion as a geometric ``memory'' rather than a new energy form. Tested against 1,590 Pantheon+ supernovae, the CRM yields $\Delta\chi^2 = -12.2$ relative to $\Lambda$CDM.

\item \textbf{Modified Newtonian Dynamics (MOND)} \cite{Milgrom1983}: MOND modifies gravitational dynamics at accelerations below $a_0 \approx 1.2 \times 10^{-10}$\,m/s$^2$, successfully predicting galactic rotation curves, the baryonic Tully-Fisher relation \cite{McGaugh2016}, and the radial acceleration relation \cite{Lelli2017} without invoking dark matter (for a comprehensive review, see \cite{FamaeyMcGaugh2012}).
\end{enumerate}

The central question of this paper is: \textit{Can both frameworks be unified into a single model that eliminates the entire dark sector?}

\subsection{The Compatibility Question}

At first glance, CRM and MOND address different ``dark'' problems:
\begin{itemize}
\item CRM replaces \textbf{dark energy} (cosmological expansion)
\item MOND replaces \textbf{dark matter} (galactic dynamics)
\end{itemize}

However, a naive combination encounters a fundamental tension: the standard CRM fits $\Omega_m \approx 0.36$, implying substantial dark matter ($\Omega_m - \Omega_b \approx 0.31$). If MOND is correct and dark matter does not exist, the model must function with $\Omega_m = \Omega_b \approx 0.05$ alone.

\subsection{Structure Formation: Common Ground}

Both frameworks converge on a critical prediction: structures form \textit{earlier and more efficiently} than $\Lambda$CDM allows.

\begin{itemize}
\item \textbf{CRM:} The later onset of cosmic acceleration ($z_{\mathrm{acc}} = 0.52$ vs.\ $0.84$) extends the matter-dominated growth phase \cite{Geiger2026}.
\item \textbf{MOND:} Enhanced gravitational attraction at low accelerations leads to faster gravitational collapse on large scales \cite{Asencio2023}.
\end{itemize}

This shared prediction is supported by multiple observational anomalies: the JWST ``Universe Breakers'' at $z > 7$ \cite{Labbe2023, BoylanKolchin2023}, the El~Gordo cluster at $z \approx 0.87$ (${>}6\sigma$ tension with $\Lambda$CDM) \cite{Asencio2023}, and unexpectedly mature protoclusters at $z > 4$ \cite{Miller2018}.


% ===================================================================
% 2. THEORETICAL FRAMEWORK
% ===================================================================
\section{Theoretical Framework}
\label{sec:theory}

\subsection{The Extended Curvature Relaxation Model}

In the standard CRM \cite{Geiger2026}, the Friedmann equation reads:
\begin{equation}
H^2(a) = H_0^2 \left[\Omega_m\,a^{-3} + \Omega_\Phi(a)\right]
\end{equation}
with
\begin{equation}
\Omega_\Phi(a) = \Phi_0 \cdot \frac{\tanh\!\big(k\cdot(a - a_{\mathrm{trans}})\big) + s}{1 + s}
\end{equation}

For the baryon-only extension, we decompose the geometric potential into two components:
\begin{equation}
\begin{split}
H^2(a) = H_0^2 \Big[&\mu_{\mathrm{eff}}\,\Omega_b\,a^{-3} + \underbrace{\Phi_0 \cdot f_{\mathrm{sat}}(a)}_{\text{geom.\ DE}} \\
&+ \underbrace{\alpha \cdot a^{-\beta}}_{\text{geom.\ DM}} + f_{\mathrm{EDE}}(a)\Big]
\end{split}
\label{eq:extended_cfm}
\end{equation}

where:
\begin{itemize}
\item $\mu_{\mathrm{eff}}\,\Omega_b$ is the MOND-enhanced baryonic matter density, with $\mu_{\mathrm{eff}} \approx 1.77$ encoding the gravitational enhancement from the MOND interpolation function at cosmological scales.\footnote{We use $\mu_{\mathrm{eff}}$ for the background-level MOND enhancement and reserve $\mu(k,a)$ for the perturbation-level modified Poisson equation (Paper~III). These are distinct quantities.}
\item $\Phi_0 \cdot f_{\mathrm{sat}}(a)$ is the saturation-type dark energy replacement (from the Dynamic Saturation Mechanism)
\item $\alpha \cdot a^{-\beta}$ is a power-law term that assumes the \textit{cosmological} role of dark matter
\item $f_{\mathrm{EDE}}(a)$ is a parametric Early Dark Energy contribution active near recombination
\end{itemize}

The flatness constraint $H^2(a{=}1)/H_0^2 = 1$ yields:
\begin{equation}
\mu_{\mathrm{eff}}\,\Omega_b + \Phi_0 \cdot f_{\mathrm{sat}}(1) + \alpha = 1
\end{equation}
Note that at $a=1$ the first term in Eq.~\eqref{eq:extended_cfm} is $\mu_{\mathrm{eff}}\,\Omega_b$, not $\Omega_b$ alone, because the MOND enhancement factor $\mu_{\mathrm{eff}}$ applies to the baryonic density at all epochs including today.

\subsection{Trace Coupling and BBN Consistency}
\label{subsec:trace_coupling}

A critical constraint on the geometric DM term is Big Bang Nucleosynthesis (BBN): at $a \sim 10^{-9}$, the naive power-law $\alpha \cdot a^{-2}$ would yield $\sim 10^{18}$, completely dominating the Friedmann equation and destroying the predicted primordial element abundances. The term \textit{must} be suppressed during the radiation era.

We propose that the geometric DM term couples not to the energy density $\rho$ but to the \textit{trace of the energy-momentum tensor}:
\begin{equation}
T \equiv g^{\mu\nu} T_{\mu\nu} = -\rho + 3p = -\rho(1 - 3w)
\end{equation}

This trace has a remarkable property: for relativistic matter (radiation, $w = 1/3$), the trace vanishes exactly:
\begin{equation}
T_{\mathrm{rad}} = -\rho_{\mathrm{rad}} + 3 \cdot \tfrac{1}{3}\rho_{\mathrm{rad}} = 0
\end{equation}

This is not a coincidence but a consequence of \textit{conformal symmetry}: massless fields are conformally invariant, and the trace of a conformally invariant energy-momentum tensor vanishes identically. During the radiation-dominated era, conformal symmetry is exact, and the geometric DM term is automatically suppressed.

For non-relativistic matter ($w \approx 0$), the trace is $T_{\mathrm{mat}} = -\rho_m \neq 0$, and the geometric DM term activates. The transition occurs naturally at matter-radiation equality ($a_{\mathrm{eq}} \approx 1.8 \times 10^{-3}$), well after BBN ($a_{\mathrm{BBN}} \sim 10^{-9}$).

The full extended Friedmann equation with trace coupling reads:
\begin{equation}
H^2(a) = H_0^2 \left[\Omega_b\,a^{-3} + \Phi_0 \cdot f_{\mathrm{sat}}(a) + \alpha \cdot a^{-\beta} \cdot \mathcal{S}(a)\right]
\label{eq:extended_cfm_trace}
\end{equation}
where $\mathcal{S}(a)$ is the trace-coupling suppression factor:
\begin{equation}
\mathcal{S}(a) = \frac{|T|}{|T| + \rho_{\mathrm{rad}}} = \frac{\Omega_b\,a^{-3}}{\Omega_b\,a^{-3} + \Omega_r\,a^{-4}}
= \frac{1}{1 + (a_{\mathrm{eq}}/a)}
\label{eq:suppression}
\end{equation}
with $a_{\mathrm{eq}} = \Omega_r/\Omega_b \approx 1.8 \times 10^{-3}$ (using $\Omega_r \approx 9 \times 10^{-5}$ and $\Omega_b \approx 0.05$ in this baryon-only model; note that if the effective matter density $\mu_{\mathrm{eff}}\,\Omega_b \approx 0.089$ is used instead, $a_{\mathrm{eq}} \approx 1.0 \times 10^{-3}$). This factor satisfies:
\begin{itemize}
\item $\mathcal{S}(a \ll a_{\mathrm{eq}}) \approx a/a_{\mathrm{eq}} \to 0$ \quad (radiation era: BBN protected)
\item $\mathcal{S}(a \gg a_{\mathrm{eq}}) \approx 1$ \quad (matter/DE era: full geometric DM)
\item $\mathcal{S}(a = 1) \approx 1 - 1.8\times10^{-3} \approx 1$ \quad (today: SN fit unchanged)
\end{itemize}

\textbf{Impact on the Pantheon+ fit:} Since all Pantheon+ supernovae are at $z < 2.3$ ($a > 0.30$), the suppression factor is $\mathcal{S} > 0.999$ throughout the observed redshift range. The MCMC results ($\alpha$, $\beta$, $\chi^2$) are unchanged to numerical precision.

\textbf{Physical interpretation:} The trace coupling has a deep geometric meaning. Physically, the geometric DM term represents the curvature ``memory'' of the initial energy concentration. During the radiation era, the universe is conformally flat (radiation is scale-free), and there is no curvature memory to sustain. The geometric DM term activates only when conformal symmetry is broken by the emergence of massive (non-relativistic) matter -- precisely at the epoch when CDM would begin to form structures in the standard picture.

\textbf{Lagrangian origin from $f(R) = R + 2\gamma R^2$:} Paper~III \cite{Geiger2026c} demonstrates that the trace-coupling suppression $\mathcal{S}(a)$ is \textit{not} an ad~hoc postulate but follows rigorously from the $R^2$ sector of the CRM Lagrangian. The trace of the field equations for $f(R) = R + 2\gamma R^2$ gives:
\begin{equation}
R + 12\gamma\,\Box R = -8\pi G\,T
\label{eq:trace_fR}
\end{equation}
where $T = g^{\mu\nu}T_{\mu\nu}$ is the matter energy-momentum trace. During the radiation era, $T_{\mathrm{rad}} = 0$ (conformal symmetry), which implies $R + 12\gamma\,\Box R = 0$. On FLRW, the decaying solution gives $R \to 0$ in the radiation era, so the $R^2$ correction -- which sources the scalaron and hence the geometric DM term -- vanishes automatically. The suppression factor $\mathcal{S}(a)$ in Eq.~\eqref{eq:suppression} is thus the phenomenological parametrization of a rigorous $f(R)$ result: the geometric DM ``activates'' precisely when $T \neq 0$ (i.e., when non-relativistic matter dominates). The only genuine postulate is the choice $f(R) = R + 2\gamma R^2$ -- the trace coupling is a \textit{consequence}, not an input. Paper~III also derives the scalaron mass $m_s \sim \mathcal{O}(10)\,H_0 \sim 10^{-32}\,\mathrm{eV}$ from the MCMC constraint $\alpha_{M,0} = 0.0011^{+0.0010}_{-0.0006}$ (48-walker chain, converged), corresponding to a cosmological Compton wavelength $\lambda_C \gtrsim 100\,\mathrm{Mpc}$ and a solar system screening length $\lambda_C^{\mathrm{solar}} \sim 20\,\mathrm{m}$ (chameleon mechanism), ensuring compatibility with local gravity tests.

\subsection{Physical Interpretation of the Geometric DM Term}

The term $\alpha \cdot a^{-\beta}$ with $\beta \approx 2.0$ (from MCMC) requires physical interpretation:

\begin{enumerate}
\item \textbf{Scaling behavior:} The MCMC posterior yields $\beta = 2.02 \pm 0.20$, consistent with curvature-like scaling ($a^{-2}$, i.e., $\beta = 2$). This is the scaling of spatial curvature in the Friedmann equation, suggesting a geometric rather than material origin.

\item \textbf{Interpretation:} In the framework of Paper~I \cite{Geiger2026}, this term represents a second equilibrium mechanism: while the saturation term describes the ``releasing brake'' (dark energy), the power-law term describes the ``geometric inertia'' of the curvature return -- a residual geometric effect that decays with expansion but slower than matter.

\item \textbf{Connection to MOND:} In the relativistic MOND theory AeST (Aether Scalar Tensor) of Skordis \& Z{\l}o\'snik \cite{Skordis2021}, a scalar field and a vector field produce an effective energy-momentum tensor that modifies the expansion history. The power-law term $\alpha \cdot a^{-\beta}$ may be interpretable as the cosmological imprint of this MOND-like modification.

\item \textbf{Effective equation of state:} The geometric DM term has an effective equation of state $w_{\mathrm{DM,geom}} = \beta/3 - 1 = -0.33 \pm 0.07$, virtually identical to the curvature equation of state ($w_k = -1/3$). The ``dark matter'' component is indistinguishable from spatial curvature.
\end{enumerate}

\subsection{MOND on Galactic vs.\ Cosmological Scales}

A key distinction must be maintained:
\begin{itemize}
\item \textbf{Galactic scales:} MOND modifies the gravitational force law below $a_0$, explaining rotation curves and the Tully-Fisher relation \textit{without dark matter}.
\item \textbf{Cosmological scales:} The extended CRM replaces dark matter's \textit{cosmological role} (contribution to $H(z)$) with a geometric potential, without requiring a particle species.
\end{itemize}

The two mechanisms are complementary: MOND handles local dynamics, while the geometric DM term handles the global expansion history.

\textit{Terminology and microscopic status.} The MOND enhancement factor $\mu_{\mathrm{eff}}$ introduced in this paper operates as an \emph{effective MOND background coupling} -- a phenomenological parametrization of how MOND-scale physics enters the Friedmann equation. Paper~IV \cite{Geiger2026d} provides the microscopic realization: a teleodynamic coupling $\mathcal{F}(A_\mu, \phi, T)$ between the scalaron, a timelike vector field $A_\mu$, and the matter trace $T$ generates the MOND enhancement dynamically, with $a_0 = cH_0/(2\pi)$ emerging as a consistency relation rather than a fitted parameter.

\subsection{The Efficiency Hypothesis: Why No Dark Matter?}
\label{subsec:efficiency}

A critical question remains: the extended CRM shows that the data \textit{permit} a universe whose matter content is exclusively baryonic, but why should this be the case? We emphasize: the CRM does not require that non-baryonic matter never existed -- only that baryons are what \textit{remains} as the matter content of the present universe, while the gravitational effects attributed to dark matter are geometric (the scalaron field). The equilibrium framework of Paper~I \cite{Geiger2026} provides a compelling answer for why baryonic matter alone is sufficient.

In the equilibrium between null space and spacetime bubble \cite{Geiger2026}, the spacetime bubble receives a finite energy budget $E_0$. Its objective is to neutralize the concentration gradient $G$ as efficiently as possible. This creates a resource allocation problem:

\begin{itemize}
\item \textbf{Baryonic matter:} Interacts electromagnetically, forms stars, produces radiation, collapses into black holes, and generates entropy at maximal rates. Baryons are \textit{highly efficient tools} for gradient reduction.

\item \textbf{Dark matter (hypothetical):} Interacts only gravitationally. It clumps but does not radiate, does not form stars, and contributes minimally to entropy production compared to an equivalent mass of baryonic matter.
\end{itemize}

In an optimally configured universe, allocating 85\% of the energy budget to a component that barely contributes to the primary objective (entropy-driven gradient reduction) would be a \textit{strategically inferior allocation}. An optimal system maximizes entropy production per unit energy by channeling the entire budget into ``active'' (baryonic) matter.

\begin{proposition}[Efficiency Principle -- Conditional Form]
\textit{If} the equilibrium between null space and spacetime bubble optimizes entropy production per unit energy (Premise~P1), and \textit{if} baryonic matter produces more entropy per unit mass than any hypothetical dark matter species (Premise~P2), \textit{then} the optimal matter content consists exclusively of baryonic matter ($\Omega_m = \Omega_b$). The gravitational effects conventionally attributed to dark matter are instead geometric consequences of the curvature return mechanism (the $\alpha \cdot a^{-\beta}$ term).
\end{proposition}

\textit{Logical structure:} The argument has the form $P1 \wedge P2 \Rightarrow S$, which is deductively valid. Premise~P2 is empirically grounded: baryons form stars, drive nucleosynthesis, and power black hole accretion, while dark matter (if it existed) would interact only gravitationally and contribute negligibly to entropy production. Premise~P1 -- that the equilibrium selects for maximal entropy production -- is the \textit{hypothesis to be tested}. It is supported by the empirical success of the model ($\Delta\chi^2 = -26.3$) but is not independently proven. The Efficiency Principle is therefore a \textit{testable conditional prediction}: it would be falsified by the experimental detection of dark matter particles.

This conditional formulation avoids circularity: we do not assume the absence of dark matter; rather, we derive it as a consequence of P1, which is itself subject to empirical test.

\textbf{Teleological objection (P1).} A persistent objection to MEPP-based and efficiency-based cosmological arguments is that they appear to attribute an \textit{objective} or \textit{goal} to the universe---``the universe optimizes entropy production''---which is a teleological structure absent from fundamental physics. The CRM framework avoids this objection to the extent that P1 is treated as a \textit{physical selection principle} in the tradition of maximum entropy statistical mechanics (Jaynes 1957) and the Principle of Least Action: the universe does not ``try'' to optimize, but configurations that happen to maximize entropy production are statistically overwhelmingly more probable (Boltzmann argument). P1 is therefore a probabilistic, not teleological, claim. However, the precise probabilistic justification for applying this principle to cosmological initial conditions has not been provided here and remains the central open problem. Until a derivation of P1 from first principles is supplied, the Efficiency Hypothesis should be treated as a heuristic organizing principle rather than a deductive argument.

The quantitative test is whether the geometric term $\alpha \cdot a^{-\beta}$ can reproduce all cosmological signatures traditionally attributed to dark matter (expansion history, CMB acoustic peaks, matter power spectrum). The Pantheon+ test presented below addresses the first of these.

\subsection{The Geometric Phase Transition}
\label{subsec:phase_transition}

The extended Friedmann equation~\eqref{eq:extended_cfm} contains two geometric terms: the power-law $\alpha \cdot a^{-\beta}$ and the saturation $\Phi_0 \cdot f_{\mathrm{sat}}(a)$. A key insight emerges: these are not independent phenomena but \textit{two phases of a single geometric process} -- the curvature return mechanism operating in different regimes.

\begin{enumerate}
\item \textbf{Early universe ($a \ll a_{\mathrm{trans}}$):} The curvature return is far from saturation. The geometric potential is dominated by the power-law term $\alpha \cdot a^{-2}$, which scales like spatial curvature and plays the cosmological role of ``dark matter'' -- providing gravitational scaffolding for structure formation.

\item \textbf{Transition epoch ($a \approx a_{\mathrm{trans}}$):} As the universe expands, the curvature return approaches its saturation limit $\Phi_0$. The power-law contribution decays, while the saturation term rises.

\item \textbf{Late universe ($a \gtrsim a_{\mathrm{trans}}$):} The saturation term dominates, providing a near-constant geometric potential that drives accelerated expansion -- the role conventionally attributed to ``dark energy.''
\end{enumerate}

This picture yields a natural interpretation: \textit{dark matter and dark energy are not two separate substances but two phases of the same geometric phenomenon.} In the early universe, spacetime geometry behaves like dark matter; in the late universe, the same geometry behaves like dark energy. The ``phase transition'' is the saturation of the curvature return potential.

The geometric potential is a decaying remnant of the Big Bang's initial curvature concentration. Early on, it provides gravitational structure (``dark matter''). As it decays and saturates, it provides accelerated expansion (``dark energy''). There is no dark sector -- only geometry at different stages of relaxation. Formally: the cosmological dark sector is a single geometric phenomenon -- the curvature return potential $\Omega_\Phi$. Its two apparent components -- dark matter ($\alpha \cdot a^{-\beta}$, dominant at early times) and dark energy ($\Phi_0 \cdot f_{\mathrm{sat}}$, dominant at late times) -- represent the unsaturated and saturated phases of the same relaxation process.


\subsection{The Running Curvature Coupling}
\label{subsec:running_beta}

The SN-only analysis with constant $\beta \approx 2.02$ yields an excellent fit to the Hubble diagram but faces a fundamental tension when confronted with CMB data: the acoustic scale $\ell_A = \pi d_C(z_*)/r_s(z_*)$ is predicted as 316.9 instead of the Planck value $301.5 \pm 0.14$, a catastrophic $>100\sigma$ discrepancy. The root cause is that $\beta = 2$ scales too slowly ($a^{-2}$) compared to CDM ($a^{-3}$), making the geometric DM term subdominant at $z \gg 1$ and yielding incorrect early-universe physics.

This tension has a natural resolution within the MOND-CRM framework. Just as MOND posits two gravitational regimes separated by the acceleration scale $a_0$:
\begin{equation}
\mu(x) = \begin{cases} 1 & \text{if } x \gg 1 \quad (\text{Newton}) \\ x & \text{if } x \ll 1 \quad (\text{MOND}) \end{cases}
\end{equation}
the curvature coupling $\beta$ should transition between two regimes separated by a curvature scale:
\begin{equation}
\beta_{\mathrm{eff}}(a) = \beta_{\mathrm{late}} + \frac{\beta_{\mathrm{early}} - \beta_{\mathrm{late}}}{1 + (a/a_t)^n}
\label{eq:running_beta}
\end{equation}
where $\beta_{\mathrm{late}} = 2.02$ (curvature-like, from the SN fit), $\beta_{\mathrm{early}} \approx 2.8$ (approaching CDM-like $a^{-3}$), $a_t$ is the transition scale factor, and $n$ controls the transition sharpness.

\textit{Physical motivation:} At early times ($z > z_t$), the Ricci scalar $R$ is enormous and the curvature return mechanism operates at full strength, mimicking CDM. As the universe expands and curvature decreases past a critical threshold, the return potential weakens and the geometric term relaxes to its low-curvature form ($a^{-2}$). The transition redshift $z_t \approx 7\text{--}10$ coincides with the epoch when the first galaxies form and MOND effects begin to manifest on galactic scales -- the same curvature threshold governs both the cosmological background transition and the onset of modified gravity on small scales.

The extended Friedmann equation with running $\beta$ and MOND background coupling reads:
\begin{equation}
\begin{split}
H^2(a) = H_0^2 \Big[&\mu_{\mathrm{eff}}\,\Omega_b\,a^{-3} + \Omega_r\,a^{-4} \\
&+ \Phi_0 \cdot f_{\mathrm{sat}}(a) + \alpha \cdot a^{-\beta_{\mathrm{eff}}(a)} \\
&+ f_{\mathrm{EDE}}(a)\Big]
\end{split}
\label{eq:running_friedmann}
\end{equation}
with $\mu_{\mathrm{eff}} = \sqrt{\pi} \approx 1.77$ (Section~\ref{subsec:sqrt_pi}) and the flatness constraint $E^2(a{=}1) = 1$ determining $\Phi_0$.


\subsection{The MOND--CRM Connection}
\label{subsec:mond_cfm_connection}

The running-$\beta$ parametrization~\eqref{eq:running_beta} reveals a deep structural connection between the CRM and MOND. Milgrom's acceleration scale $a_0 \approx 1.2 \times 10^{-10}$\,m/s$^2$ is empirically related to the Hubble parameter via:
\begin{equation}
a_0 \approx \frac{c\,H_0}{2\pi} \approx \frac{c\,H_0}{5}
\end{equation}
In the CRM framework, this coincidence acquires a natural explanation: the Hubble acceleration $a_H = Hc$ sets the curvature scale, and the MOND threshold $a_0$ marks the boundary between the high-curvature (CDM-like) and low-curvature (geometric) regimes. The three phases of the CRM+MOND universe are:

\begin{enumerate}
\item \textbf{Phase~1 ($z > z_t \approx 7$, Newton/CDM regime):} $\beta_{\mathrm{eff}} \approx 2.8$, strong curvature return provides deep potential wells. On galactic scales, accelerations exceed $a_0$. The universe behaves identically to $\Lambda$CDM.

\item \textbf{Phase~2 ($z \sim z_t$, transition):} $\beta_{\mathrm{eff}}$ decreases from $\sim$3 to $\sim$2. Curvature return weakens, gravitational potentials flatten. On galactic scales, accelerations approach $a_0$ and MOND effects begin to manifest.

\item \textbf{Phase~3 ($z < z_t$, MOND/geometric regime):} $\beta_{\mathrm{eff}} \approx 2.02$, the geometric term scales as curvature ($a^{-2}$). On galactic scales, $a < a_0$ and MOND produces flat rotation curves without CDM halos. Cosmologically, the saturation term drives acceleration.
\end{enumerate}

This three-phase structure provides a causal chain: \textit{the end of curvature return triggers both the cosmological acceleration and the galactic MOND regime} -- they are two manifestations of the same geometric transition.


\subsection{The $\sqrt{\pi}$ Conjecture: Dimensional Origin of $\mu_{\mathrm{eff}}$}
\label{subsec:sqrt_pi}

The refined joint fit (Section~\ref{subsec:joint_fit}) yields $\mu_{\mathrm{eff}} = 1.77$, which is numerically indistinguishable from $\sqrt{\pi} = 1.7725$ (deviation 0.2\%). We propose -- as a conjecture requiring independent derivation -- that this may reflect the \textit{dimensional geometry} of the gravitational phase space rather than a numerical coincidence.

\textbf{Key observation:} The volumes of the unit $n$-sphere are:
\begin{equation}
V_1 = 2, \qquad V_2 = \pi, \qquad V_3 = \frac{4\pi}{3}
\end{equation}
The standard MOND enhancement factor for galactic rotation curves is:
\begin{equation}
\mu_{\mathrm{eff}}^{\mathrm{gal}} = \frac{V_3}{V_2} = \frac{4}{3}
\label{eq:mu_gal}
\end{equation}
-- the ratio of the 3D (sphere) to 2D (disk) gravitational phase space volumes. This is the classic MOND $4/3$ factor, appropriate for galaxies where gravity operates in 3D and the disk geometry of the galaxy projects onto a 2D plane.

For cosmological expansion, the geometry is fundamentally different: the Friedmann equation describes a homogeneous, isotropic universe where the expansion is parametrized by a single scale factor $a(t)$. The relevant geometric quantity is not a ratio of phase space volumes but the \textit{projection amplitude} of the 2-sphere onto the observational space:
\begin{equation}
\mu_{\mathrm{eff}}^{\mathrm{cosmo}} = \sqrt{V_2} = \sqrt{\pi} \approx 1.7725
\label{eq:mu_cosmo}
\end{equation}

\begin{conjecture}[$\sqrt{\pi}$ Conjecture]
The MOND gravitational enhancement at cosmological scales is $\mu_{\mathrm{eff}} = \sqrt{\pi}$, arising from the square root of the unit-circle area. This connects to the Gaussian integral $\Gamma(1/2) = \sqrt{\pi} = \int_{-\infty}^{\infty} e^{-t^2}\,dt$ and reflects the thermodynamic normalization of gravitational modes on the cosmological 2-sphere.
\end{conjecture}

\textbf{Numerological caveat.} The agreement between the fitted $\mu_{\mathrm{eff}} = 1.77$ and $\sqrt{\pi} \approx 1.7725$ at the 0.2\% level is suggestive but does not constitute a derivation. Numerical coincidences of this precision are not uncommon in physics without physical content (Numerology problem). The conjecture is physically motivated by the geometric argument above, but this motivation is post-hoc: the geometric interpretation was found \emph{after} fitting the data to the numerical value, not derived \emph{before} fitting. An independent theoretical derivation of $\sqrt{\pi}$ from the MOND action or the CRM Lagrangian is required to distinguish a physically meaningful connection from a numerical accident. Until such a derivation exists, the $\sqrt{\pi}$ Conjecture should not be used in the abstract or introduction as if it were established physics.

This conjecture yields remarkable quantitative predictions. Setting $\mu_{\mathrm{eff}} = \sqrt{\pi}$ \textit{exactly} and fitting the remaining 6 parameters against SN + CMB + BAO data yields:
\begin{center}
\begin{tabular}{lcc}
\toprule
& CRM+MOND ($\mu = \sqrt{\pi}$) & $\Lambda$CDM \\
\midrule
$H_0$ [km/s/Mpc] & $\mathbf{69.0}$ & 67.4 \\
$r_d$ [Mpc] & 143.3 & 147.2 \\
$\ell_A$ & 301.471 & 301.428 \\
$\mathcal{R}$ & 1.7502 & 1.7496 \\
$\chi^2_{\mathrm{total}}$ & $\mathbf{704.2}$ & 710.3 \\
$\Delta\chi^2$ & $\mathbf{-6.1}$ & 0 \\
\bottomrule
\end{tabular}
\end{center}
With $H_0 = 69$\,km/s/Mpc, the model sits \textit{between} the Planck ($67.4$) and SH0ES ($73.0$) values, potentially bridging the Hubble tension. The effective baryon density is:
\begin{equation}
\Omega_{b,\mathrm{eff}} = \sqrt{\pi}\,\Omega_b \approx 1.77 \times 0.047 = 0.083
\end{equation}
and the total ``matter-like'' contribution is:
\begin{equation}
3\sqrt{\pi}\,\Omega_b \approx 0.250 \approx \Omega_{\mathrm{CDM}}^{\Lambda\mathrm{CDM}}
\end{equation}
This suggests that the dark matter density in $\Lambda$CDM ($\Omega_{\mathrm{CDM}} \approx 0.265$) is a geometric artifact: $3\sqrt{\pi}$ times the physical baryon density. With updated Planck~2018 best-fit values ($\Omega_b = 0.0493$, $\Omega_{\mathrm{CDM}} = 0.2660$), the prediction $3\sqrt{\pi}\,\Omega_b = 0.2606$ matches the observed $\Omega_{\mathrm{CDM}}$ to \textbf{1.3\%} -- a falsifiable, parameter-free prediction.

\textit{Connection to the saturation ODE:} The factor $\sqrt{\pi}$ appears naturally in the thermodynamics of the CRM. The saturation ODE $d\Omega_\Phi/da = k(1 - \Omega_\Phi^2/\Phi_0^2)$ has a $\tanh$ solution, and the partition function of this $\cosh^{-2}$ (P\"oschl-Teller) system involves Gaussian integrals with $\sqrt{\pi}$ normalization. The Gel'fand-Yaglom method applied to the P\"oschl-Teller operator on $[0,L]$ yields functional determinant ratios $\sqrt{\det} = 1.677$ for $\lambda = 1$, close to $\sqrt{\pi} = 1.7725$. The same mathematical structure that generates the saturation mechanism also determines the MOND enhancement factor on the cosmological background.


% ===================================================================
% 3. DATA ANALYSIS
% ===================================================================
\section{Data Analysis and Results}
\label{sec:analysis}

\subsection{Data and Methodology}

We use the Pantheon+ catalog \cite{Scolnic2022} comprising 1,590 Type~Ia supernovae with $z > 0.01$ (redshift range $0.01$--$2.26$). Luminosity distances are computed via cumulative trapezoidal integration on a fine redshift grid ($N = 2{,}000$). The nuisance parameter~$M$ is analytically marginalized. Parameter optimization uses differential evolution with L-BFGS-B polish.

\subsection{Results: Model Comparison}

\begin{table*}
\centering
\caption{Model comparison against 1,590 Pantheon+ supernovae.}
\label{tab:comparison}
\begin{tabular}{lcccccc}
\toprule
\textbf{Model} & $\Omega_m$ & \textbf{Params} & $\chi^2$ & $\Delta\chi^2$ & AIC & BIC \\
\midrule
$\Lambda$CDM & 0.244 & 2 & 729.0 & 0 & 733.0 & 743.7 \\
CRM Standard & 0.364 & 4 & 716.8 & $-12.2$ & 724.8 & 746.3 \\
\midrule
CRM Baryon Fixed & 0.050 & 3 & 945.5 & $+216.5$ & 951.5 & 967.6 \\
CRM Baryon Band & 0.070 & 4 & 894.7 & $+165.7$ & 902.7 & 924.1 \\
\midrule
\textbf{Extended CRM+MOND} & \textbf{0.050} & \textbf{5} & \textbf{702.7} & $\mathbf{-26.3}$ & \textbf{712.7} & 739.5 \\
\bottomrule
\end{tabular}
\end{table*}

\subsection{Key Findings}

\begin{enumerate}
\item \textbf{Simple baryon-only CRM fails:} With $\Omega_m = 0.05$ and only the $\tanh$ saturation term, the fit degrades catastrophically ($\Delta\chi^2 = +216.5$). The optimizer attempts extreme parameters ($k = 86$, $a_{\mathrm{trans}} = 0.06$) to create a near-step-function, confirming that the standard CRM \textit{cannot} compensate for missing dark matter.

\item \textbf{Extended CRM achieves improved fit:} Adding the geometric DM term $\alpha \cdot a^{-\beta}$ restores and exceeds the fit quality, achieving $\Delta\chi^2 = -26.3$ versus $\Lambda$CDM across the Pantheon+ SN dataset.

\item \textbf{Best-fit parameters (MCMC):} A full Markov Chain Monte Carlo analysis (emcee, 48 walkers, 5000 steps) yields:
\begin{itemize}
\item Saturation term: $\Phi_0 = 0.43^{+0.14}_{-0.08}$, $k = 9.8^{+6.7}_{-3.8}$, $a_{\mathrm{trans}} = 0.971^{+0.016}_{-0.031}$ ($z_{\mathrm{trans}} = 0.03$)
\item Geometric DM term: $\alpha = 0.68^{+0.02}_{-0.07}$, $\beta = 2.02^{+0.26}_{-0.14}$
\item Energy budget at $a = 1$: $\Omega_b = 0.05$, $\Omega_\Phi = 0.95$ (total geometric contribution)
\end{itemize}

\item \textbf{$\beta \approx 2.0$: Curvature scaling.} The MCMC posterior for $\beta$ peaks at $2.02 \pm 0.20$, consistent with \textit{curvature-like scaling} ($a^{-2}$, i.e., $w = -1/3$). This is a remarkable result: the data independently recover a scaling exponent that corresponds to \textit{spatial curvature}, not to a material component. The effective equation of state $w_{\mathrm{DM,geom}} = \beta/3 - 1 = -0.33$ is virtually identical to the curvature equation of state.

\item \textbf{Late saturation transition:} The saturation transition occurs very late ($z_{\mathrm{trans}} \approx 0.03$), much later than in the standard CRM ($z_{\mathrm{trans}} = 0.33$). The geometric DM term (curvature-like) dominates the early expansion, while the saturation term provides the late-time acceleration.

\item \textbf{AIC vs.\ BIC:} The $\Delta\mathrm{AIC} = -16.3$ strongly favors the extended model. The $\Delta\mathrm{BIC} = -4.2$ also favors it despite the parameter penalty (5 vs.\ 2 parameters). This is the first model in our analysis to achieve \textit{both} AIC and BIC preference over $\Lambda$CDM simultaneously.
\end{enumerate}


\subsection{Cross-Validation: Ruling Out Overfitting}
\label{subsec:crossval}

With five free parameters versus two for $\Lambda$CDM, an overfitting concern is natural. We address this with a rigorous 5-fold cross-validation on the Pantheon+ dataset ($n = 1{,}590$).

\textbf{Method:} The data are randomly split into five equal folds (seed $= 42$). For each fold, both models ($\Lambda$CDM and Extended CRM) are optimized on the remaining 80\% (training set) using differential evolution, and the predictive $\chi^2/n$ is evaluated on the held-out 20\% (test set). This procedure tests \textit{generalization}, not merely goodness-of-fit.

\begin{table*}
\centering
\caption{5-fold cross-validation results: predictive $\chi^2/n$ on held-out test sets.}
\label{tab:crossval}
\begin{tabular}{lcccccc}
\toprule
\textbf{Model} & \textbf{Fold 1} & \textbf{Fold 2} & \textbf{Fold 3} & \textbf{Fold 4} & \textbf{Fold 5} & $\langle\chi^2/n\rangle$ \\
\midrule
$\Lambda$CDM (2 params) & 0.467 & 0.428 & 0.461 & 0.435 & 0.468 & $0.452 \pm 0.017$ \\
Ext.\ CRM+MOND (5 params) & 0.456 & 0.428 & 0.465 & 0.411 & 0.467 & $0.445 \pm 0.022$ \\
\bottomrule
\end{tabular}
\end{table*}

\textbf{Result:} The Extended CRM achieves a \textit{lower} mean predictive $\chi^2/n$ on unseen data ($\Delta\langle\chi^2/n\rangle = -0.007$). Despite having 2.5$\times$ more parameters, the model generalizes \textit{better} than $\Lambda$CDM, ruling out overfitting as an explanation for the $\Delta\chi^2 = -26.3$ improvement. The slightly higher fold-to-fold variance ($\sigma = 0.022$ vs.\ $0.017$) is expected for a more flexible model but does not indicate instability.


\subsection{Joint SN + CMB + BAO Fit with Running $\beta$}
\label{subsec:joint_fit}

The SN-only analysis demonstrates the viability of the baryon-only framework at $z \lesssim 2.3$. The decisive test, however, is joint compatibility with CMB and BAO data. We use three data sets simultaneously:

\begin{enumerate}
\item \textbf{Pantheon+ SN:} 1,590 Type~Ia supernovae ($z > 0.01$), analytically marginalized over $M$.
\item \textbf{Planck CMB (compressed):} Acoustic scale $\ell_A = 301.471 \pm 0.14$ and shift parameter $\mathcal{R} = 1.7502 \pm 0.0046$ at $z_* = 1089.80$.
\item \textbf{BAO:} 9 distance measurements from 6dFGS ($z = 0.15$), BOSS DR12 ($z = 0.38, 0.51, 0.61$), and Lyman-$\alpha$ ($z = 2.334$).
\end{enumerate}

\subsubsection{The CMB Catastrophe of Constant $\beta$}

With constant $\beta = 2.02$, the model produces $\ell_A = 316.9$ and $\mathcal{R} = 1.00$ -- both catastrophically wrong (Planck: $301.5$ and $1.75$). The sound horizon is $r_d = 200$\,Mpc instead of 147\,Mpc. The total $\Delta\chi^2 = +2{,}630$, ruling out the constant-$\beta$ model at $>50\sigma$.

\subsubsection{Running $\beta$ Alone}

Introducing the running coupling~\eqref{eq:running_beta} dramatically improves the fit. A grid scan over $(\beta_{\mathrm{early}}, a_t, \alpha, H_0)$ followed by Nelder-Mead optimization yields the point estimates below. \textit{Note on optimization methodology:} The grid scan + Nelder-Mead approach identifies a best-fit point but does not provide formal parameter uncertainties or posterior distributions. A full Markov Chain Monte Carlo (MCMC) analysis would be required to assess parameter degeneracies and confidence intervals; the results below should therefore be regarded as preliminary point estimates subject to this caveat.

\begin{table}[H]
\centering
\caption{Running-$\beta$ CRM: optimized parameters and comparison with $\Lambda$CDM.}
\label{tab:running_beta}
\begin{tabular}{lccc}
\toprule
\textbf{Parameter} & \textbf{CRM (running $\beta$)} & \textbf{$\Lambda$CDM} & \textbf{Planck} \\
\midrule
$\beta_{\mathrm{early}}$ & 2.823 & -- & -- \\
$\beta_{\mathrm{late}}$ & 2.02 (fixed) & -- & -- \\
$a_t$ ($z_t$) & 0.0924 (9.8) & -- & -- \\
$\alpha$ & 0.628 & -- & -- \\
$H_0$ [km/s/Mpc] & 60.0 & 67.4 & 67.4 \\
\midrule
$\ell_A$ & 301.42 & 301.43 & $301.47 \pm 0.14$ \\
$\mathcal{R}$ & 1.759 & 1.750 & $1.750 \pm 0.005$ \\
$r_d$ [Mpc] & 179.3 & 147.2 & $147.2$ \\
\midrule
$\chi^2_{\mathrm{SN}}$ & 743.3 & 700.9 & -- \\
$\chi^2_{\mathrm{CMB}}$ & 3.8 & 0.1 & -- \\
$\chi^2_{\mathrm{BAO}}$ & 33.6 & 9.3 & -- \\
$\chi^2_{\mathrm{total}}$ & 780.6 & 710.3 & -- \\
$\Delta\chi^2$ & $+70.3$ & 0 & -- \\
\bottomrule
\end{tabular}
\end{table}

The running-$\beta$ model reduces $\Delta\chi^2$ from $+2{,}630$ to $+70.3$ -- a factor-of-37 improvement. The CMB observables $\ell_A$ and $\mathcal{R}$ are now within $1\sigma$ of Planck. The remaining tension comes primarily from the sound horizon ($r_d = 179$ vs.\ $147$\,Mpc), which affects the BAO fit.

\subsubsection{Combined Fit: Running $\beta$ + Early Dark Energy}

The residual $r_d$ tension is addressed by a parametric Early Dark Energy (EDE) contribution active near recombination:
\begin{equation}
f_{\mathrm{EDE}}(a) = \frac{f_{\mathrm{amp}}}{1 + (a/a_{\mathrm{EDE}})^p} - f_{\mathrm{EDE}}(1)
\end{equation}
which reduces the sound horizon by increasing $H(z)$ during the drag epoch without affecting low-$z$ observables. In the CRM framework, this term has a natural interpretation as residual curvature energy that has not yet converted to radiation at the epoch of recombination.

The combined 6-parameter fit ($\beta_{\mathrm{early}}$, $a_t$, $\alpha$, $H_0$, $f_{\mathrm{amp}}$, $a_{\mathrm{EDE}}$) yields the following joint fit result:

\begin{table}[H]
\centering
\caption{Preferred CRM+MOND vs.\ $\Lambda$CDM: joint SN + CMB + BAO fit. The CRM uses scale-dependent $\mu(a)$ with zero EDE and 6 free parameters -- identical to $\Lambda$CDM.}
\label{tab:combined_fit}
\begin{tabular}{lcc}
\toprule
\textbf{Observable} & \textbf{CRM+MOND} & \textbf{$\Lambda$CDM} \\
\midrule
$\mu_{\mathrm{eff}}(z{=}0)$ & $\sqrt{\pi}$ & -- \\
$f_{\mathrm{EDE}}(z_*)$ & $0\%$ & -- \\
$H_0$ [km/s/Mpc] & $67.3$ & 67.4 \\
$r_d$ [Mpc] & $146.9$ & 147.2 \\
Parameters & $6$ & 6 \\
\midrule
$\ell_A$ & $301.471$ & 301.428 \\
$\mathcal{R}$ & $1.7502$ & 1.7496 \\
\midrule
$\chi^2_{\mathrm{SN}}$ & $699.6$ & 700.9 \\
$\chi^2_{\mathrm{CMB}}$ & $0.0$ & 0.1 \\
$\chi^2_{\mathrm{BAO}}$ & $5.2$ & 9.3 \\
\midrule
$\chi^2_{\mathrm{total}}$ & $704.8$ & 710.3 \\
$\Delta\chi^2$ & $-5.5$ & 0 \\
\bottomrule
\end{tabular}
\end{table}

\textbf{Key results:}
\begin{enumerate}
\item \textbf{EDE eliminated:} The preferred variant uses scale-dependent $\mu(a)$ (Section~\ref{subsubsec:mu_scale}) that transitions from $\mu = \sqrt{\pi}$ today to $\mu \to 1$ at $z > 4000$. This \textbf{completely eliminates} the EDE component ($f_{\mathrm{EDE}} = 0$), reducing the model to 6 free parameters -- the same count as $\Lambda$CDM.

\item \textbf{Hubble constant resolved:} With scale-dependent $\mu(a)$, $H_0 = 67.3$\,km/s/Mpc -- within $0.2$\,km/s/Mpc of the Planck value ($67.4 \pm 0.5$). The constant-$\mu$ variants can reach any $H_0$ between 58 and 85\,km/s/Mpc by adjusting $\mu_{\mathrm{eff}}$.

\item \textbf{Sound horizon solved:} The sound horizon $r_d = 146.9$\,Mpc is \textit{essentially identical} to $\Lambda$CDM ($147.2$\,Mpc, deviation 0.2\%), compared to the previous 12\% discrepancy ($r_d = 165$\,Mpc).

\item \textbf{CMB:} The acoustic scale $\ell_A = 301.471$ and the shift parameter $\mathcal{R} = 1.7502$ match Planck \textit{exactly}.

\item \textbf{Overall:} $\Delta\chi^2 = -5.5$ (preferred $\mu(a)$ variant, no EDE; $\chi^2_{\mathrm{total}} = 704.8$) vs.\ $\Delta\chi^2 = -6.1$ (constant $\mu = \sqrt{\pi}$ with EDE; $\chi^2_{\mathrm{total}} = 704.2$). These are two \emph{different model variants}; the different chi-squared values reflect different physics, not an inconsistency. The preferred variant ($\Delta\chi^2 = -5.5$, no EDE) uses 6 free parameters identical in number to $\Lambda$CDM. The constant-$\mu$ variant requires $f_{\mathrm{EDE}} = 59\%$, which is not independently motivated.
\end{enumerate}

The best-fit parameters (preferred $\mu(a)$ variant) are: $\mu_{\mathrm{late}} = \sqrt{\pi}$, $\mu_{\mathrm{early}} = 1.00$, $a_\mu = 2.55 \times 10^{-4}$ ($z_\mu = 3918$), $\beta_{\mathrm{early}} = 2.82$, $a_t = 0.098$ ($z_t = 9.2$), $\alpha = 0.695$, $H_0 = 67.3$\,km/s/Mpc, $f_{\mathrm{EDE}} = 0$.

\subsubsection{The MOND Background Coupling $\mu_{\mathrm{eff}}$}
\label{subsubsec:mu_eff}

The central innovation of the refined fit is applying the MOND gravitational enhancement on the \textit{background level} of the Friedmann equation:
\begin{equation}
\begin{split}
H^2(a) = H_0^2 \Big[&\mu_{\mathrm{eff}}\,\Omega_b\,a^{-3} + \Omega_r\,a^{-4} \\
&+ \Phi_0 \cdot f_{\mathrm{sat}}(a) + \alpha \cdot a^{-\beta_{\mathrm{eff}}(a)} \\
&+ f_{\mathrm{EDE}}(a)\Big]
\end{split}
\label{eq:mond_friedmann}
\end{equation}
with the modified sound horizon:
\begin{align}
r_d &= \int_0^{a_d} \frac{c_s(a)}{a^2 H(a)} \, da\,,\quad c_s = \frac{1}{\sqrt{3(1 + R_b a)}}\,,\notag\\
R_b &= \frac{3\,\mu_{\mathrm{eff}}\,\Omega_b}{4\,\Omega_\gamma}
\label{eq:mond_rs}
\end{align}

A $\mu_{\mathrm{eff}}$-profile analysis (optimizing all other parameters at each fixed $\mu_{\mathrm{eff}}$) reveals a clear $\mu$--$H_0$ relationship:

\begin{center}
\begin{tabular}{ccc}
\toprule
$\mu_{\mathrm{eff}}$ & $H_0$ [km/s/Mpc] & $r_d$ [Mpc] \\
\midrule
1.00 & 60.0 & 165 \\
1.33 & 58.2 & 170 \\
1.50 & 64.0 & 155 \\
1.60 & 71.4 & 139 \\
1.77 & 66.0 & 150 \\
2.00 & 84.5 & 117 \\
\bottomrule
\end{tabular}
\end{center}

The non-monotonic behavior (dip at $\mu = 1.33$, then sharp rise) reflects the interplay between the MOND coupling and the EDE fraction: higher $\mu_{\mathrm{eff}}$ reduces $r_d$, which relaxes the $\ell_A$ constraint and permits higher $H_0$.

\subsubsection{Scale-Dependent $\mu(a)$: Eliminating EDE}
\label{subsubsec:mu_scale}

The constant-$\mu_{\mathrm{eff}}$ variants all require a large EDE fraction ($f_{\mathrm{EDE}} \approx 50$--$60\%$). This can be entirely eliminated by promoting $\mu$ to a \textit{scale-dependent} function:
\begin{equation}
\mu(a) = \mu_{\mathrm{late}} + \frac{\mu_{\mathrm{early}} - \mu_{\mathrm{late}}}{1 + (a/a_\mu)^4}
\label{eq:mu_running}
\end{equation}
with $\mu_{\mathrm{late}} = \sqrt{\pi}$ (the $\sqrt{\pi}$ Conjecture value, Section~\ref{subsec:sqrt_pi}) and $\mu_{\mathrm{early}}$, $a_\mu$ as free parameters. This parametrization is the \textit{exact analogue} of the running $\beta$ (Eq.~\ref{eq:running_beta}): just as $\beta$ transitions from CDM-like to curvature-like behavior, $\mu$ transitions from standard gravity ($\mu = 1$, no MOND enhancement) at very early times to full MOND enhancement ($\mu = \sqrt{\pi}$) at late times.

\textbf{Physical motivation:} At $z > 4000$, the cosmological acceleration is far above the MOND scale $a_0$, so standard gravity applies ($\mu \to 1$). As the universe expands and the effective acceleration drops below $a_0$, the MOND enhancement gradually activates ($\mu \to \sqrt{\pi}$). This is precisely the cosmological analogue of the galactic MOND transition.

The fit with scale-dependent $\mu(a)$ and \textbf{zero EDE} yields:
\begin{center}
\begin{tabular}{lcc}
\toprule
& CRM+MOND ($\mu(a)$, no EDE) & $\Lambda$CDM \\
\midrule
$H_0$ [km/s/Mpc] & $\mathbf{67.3}$ & 67.4 \\
$r_d$ [Mpc] & $\mathbf{146.9}$ & 147.2 \\
$\ell_A$ & 301.471 & 301.428 \\
$\mathcal{R}$ & 1.7502 & 1.7496 \\
$\chi^2_{\mathrm{total}}$ & $\mathbf{704.8}$ & 710.3 \\
$\Delta\chi^2$ & $\mathbf{-5.5}$ & 0 \\
Parameters & 6 & 6 \\
\bottomrule
\end{tabular}
\end{center}

The $\mu(a)$-profile at different redshifts:
\begin{center}
\begin{tabular}{rc}
\toprule
$z$ & $\mu(z)$ \\
\midrule
0 & 1.772 \\
1 & 1.772 \\
10 & 1.772 \\
100 & 1.772 \\
500 & 1.772 \\
1090 & 1.768 \\
5000 & 1.212 \\
\bottomrule
\end{tabular}
\end{center}

\textbf{Key observations:}
\begin{enumerate}
\item The MOND enhancement is essentially constant ($\mu \approx \sqrt{\pi}$) from $z = 0$ to $z \approx 1000$, only deviating at $z > 3000$.
\item The transition redshift $z_\mu \approx 3918$ is \textit{between} recombination ($z_* = 1090$) and the matter-radiation equality epoch ($z_{\mathrm{eq}} \approx 3400$).
\item The 6 free parameters ($\beta_{\mathrm{early}}$, $a_t$, $\alpha$, $H_0$, $\mu_{\mathrm{early}}$, $a_\mu$) are the \textit{same count} as $\Lambda$CDM ($\Omega_b h^2$, $\Omega_c h^2$, $H_0$, $\tau$, $n_s$, $A_s$).
\item $H_0 = 67.3$\,km/s/Mpc and $r_d = 146.9$\,Mpc are \textit{essentially identical} to the $\Lambda$CDM values.
\end{enumerate}

This result represents the most parsimonious version of the CRM+MOND framework: it eliminates not only the dark sector but also the ad-hoc EDE component, reducing the model to a minimal set of geometric parameters.

\subsubsection{Remaining Caveats}
\label{subsubsec:caveats}

Two aspects require further investigation:
\begin{enumerate}
\item \textbf{Perturbation theory:} The background observables ($\ell_A$, $\mathcal{R}$, $r_d$, $H_0$) are fully matched. Using hi\_class \cite{Zumalacarregui2017} with $f(R)$-type perturbation modifications ($\alpha_M = 0.0007$, $\alpha_B = -\alpha_M/2$, $\alpha_T = 0$), we obtain -- directly verified:
\begin{itemize}
\item $\ell_1 = 220$ -- \textbf{exact Planck match} (via early ISW effect)
\item $\mathcal{P}_3/\mathcal{P}_1 = 0.4295$ -- \textbf{exact Planck match} (via $\beta_{\mathrm{early}} = 2.83$)
\end{itemize}
BBN is fully consistent ($\Delta N_{\mathrm{eff}} \approx 0$). The angular acoustic scale offset ($100\,\theta_s = 1.034$ vs.\ Planck's $1.041$, $0.69\%$) requires careful assessment: while $0.69\%$ appears small, Planck's measurement precision is $\sigma(\theta_s) \approx 0.0003$, making this offset $>20\sigma$ -- a formally severe tension. The physical origin of the offset is clear and lies within this paper's phenomenological framework: the geometric power-law term uses $w_{\mathrm{eff}} = (\beta-3)/3 \approx -0.06$ at recombination rather than $w = 0$ (pressureless). This 6\% equation-of-state deviation causes a $\sim$2.5\% enlargement of the sound horizon ($r_s = 148.1$\,Mpc vs.\ $144.5$\,Mpc for $\Lambda$CDM), which shifts $\theta_s$ by the observed 0.69\%. The resolution within the CRM framework can be stated without delegating to Paper~III: a proper Lagrangian treatment of the $R^2$ scalaron (which constitutes the geometric dark matter) gives a CDM-like equation of state $w \approx 0$, because the scalaron is a massive scalar field whose kinetic energy is negligible during matter domination. The phenomenological power-law ansatz $\alpha a^{-\beta}$ approximates this behavior but introduces a spurious pressure term. This paper therefore has a known, physically understood source of $\theta_s$ tension that would be resolved by using the full scalaron dynamics instead of the power-law parametrization.

\item \textbf{Lagrangian derivation:} Both the running $\beta_{\mathrm{eff}}(a)$ and the running $\mu(a)$ are phenomenological transition functions. A derivation from an action principle -- potentially connecting the $\sqrt{\pi}$ conjecture to the CRM saturation mechanism -- is required for theoretical completeness (Paper~III).
\end{enumerate}


% ===================================================================
% 4. DISCUSSION
% ===================================================================
\section{Discussion}
\label{sec:discussion}

\subsection{A Universe Without a Dark Sector}

The extended CRM demonstrates that the entire expansion history probed by Type~Ia supernovae can be described with:
\begin{itemize}
\item Baryonic matter ($\Omega_b = 0.05$) -- the \textit{only} material content
\item A saturation-type geometric potential -- replacing dark energy
\item A power-law geometric term -- replacing dark matter's cosmological role
\end{itemize}

If this result survives tests against CMB and BAO data, it would imply that 95\% of the $\Lambda$CDM energy budget is an artifact of interpreting geometric effects as material components.

\subsection{Status of the MOND Coupling (Provisional)}
\label{subsec:mond_status}

The MOND background coupling $\mu_{\mathrm{eff}}$ used in this paper is a \emph{phenomenological parametrization}: its value ($\mu_{\mathrm{eff}} = \sqrt{\pi}$) and its scale dependence ($\mu(a) \to 1$ at $z > 4000$) are fitted to the data, not derived from a Lagrangian. This is an honest limitation of the present work: the coupling reproduces the correct expansion history, but its physical origin remains open.

Paper~IV \cite{Geiger2026d} provides the dynamical foundation. The vector sector -- a timelike unit field $A_\mu$ derived from the Maximum Entropy Production Principle -- generates the MOND enhancement through a teleodynamic coupling to the scalaron. The separation of concerns is:
\begin{itemize}
\item \textbf{This paper (Paper~II):} Demonstrates \emph{that} the MOND background coupling achieves joint SN+CMB+BAO compatibility and determines the effective magnitude $\mu_{\mathrm{eff}} \approx \sqrt{\pi}$.
\item \textbf{Paper~IV:} Demonstrates \emph{why} -- deriving the coupling from field equations and establishing $a_0 = cH_0/(2\pi)$ as a consistency relation.
\end{itemize}

\subsection{The $\beta \approx 2.0$ Result: Curvature as Dark Matter}

The MCMC posterior for the scaling exponent yields $\beta = 2.02^{+0.26}_{-0.14}$, remarkably close to -- and statistically consistent with -- the curvature scaling $\beta = 2$ ($a^{-2}$). This corresponds to an effective equation of state $w_{\mathrm{DM,geom}} = -0.33$, indistinguishable from spatial curvature ($w_k = -1/3$). For comparison, the standard cosmological components scale as:
\begin{itemize}
\item Matter: $\beta = 3$ ($a^{-3}$, $w = 0$)
\item Curvature: $\beta = 2$ ($a^{-2}$, $w = -1/3$) $\quad\leftarrow$ \textbf{recovered by MCMC}
\item Radiation: $\beta = 4$ ($a^{-4}$, $w = 1/3$)
\end{itemize}

This result has profound implications: the component traditionally identified as ``dark matter'' in the Friedmann equation may in fact be \textit{spatial curvature} -- not the global curvature $k$ of the FLRW metric, but a \textit{dynamic, decaying curvature memory} encoded in the geometric potential. In the framework of Paper~I, this is precisely the ``geometric inertia'' of the curvature return: a residual imprint of the Big Bang's energy concentration that dilutes with expansion at the curvature rate $a^{-2}$ rather than the matter rate $a^{-3}$.

\subsection{Relation to AeST and Relativistic MOND}

The relativistic MOND theory AeST (Aether Scalar Tensor) \cite{Skordis2021}, which supersedes Bekenstein's earlier TeVeS theory \cite{Bekenstein2004} by achieving compatibility with gravitational wave speed constraints and the CMB, provides the only known framework that simultaneously:
\begin{enumerate}
\item Reproduces MOND dynamics on galactic scales
\item Fits the CMB power spectrum (including the third acoustic peak)
\item Fits the matter power spectrum
\end{enumerate}

AeST achieves this through a scalar field $\phi$ and a timelike vector field $A_\mu$ that produce an effective energy-momentum tensor. The cosmological background equations in AeST contain terms that contribute to $H^2(a)$ with non-standard scaling. A detailed comparison between the AeST background equations and the extended CRM Friedmann equation~\eqref{eq:extended_cfm} is a key objective for future work.

\subsection{Addressing the Ultimate Cosmological Challenges}
\label{subsec:endgegner}

Any theory that eliminates dark matter must confront three critical observational pillars of $\Lambda$CDM. We address each in turn, showing how the geometric dark sector replacement provides a pathway through each challenge.

\subsubsection{Challenge 1: CMB Acoustic Peaks}

The relative heights of the CMB acoustic peaks -- particularly the ratio of the first to the third peak -- are conventionally interpreted as evidence for a gravitational component that does not interact with photons (i.e., dark matter). In $\Lambda$CDM, cold dark matter provides gravitational potential wells that drive baryon-photon oscillations without experiencing radiation pressure.

\textit{Quantitative assessment:} With constant $\beta = 2.02$, the geometric DM term contributes only $\sim$0.5\% of the baryonic density at $z_*$, because $a^{-2}$ scales slower than matter ($a^{-3}$). However, the running-$\beta$ extension (Section~\ref{subsec:running_beta}) resolves this: with $\beta_{\mathrm{eff}}(z_* = 1090) \approx 2.78$, the geometric term scales nearly as $a^{-2.8}$ -- close to CDM ($a^{-3}$) -- and contributes substantially at recombination. The effective matter density $\Omega_{m,\mathrm{eff}}$ at $z_* \approx 1090$ is 0.57 in the combined fit (Table~\ref{tab:combined_fit}), providing sufficient gravitational depth for baryon-photon oscillations. At intermediate redshifts ($z \sim 2$), $\Omega_{m,\mathrm{eff}} \approx 0.31$ -- nearly identical to $\Lambda$CDM's $\Omega_m = 0.315$.

The \textit{background-level} CMB observables are now fully reproduced: $\ell_A = 301.477$ and $\mathcal{R} = 1.7502$ match Planck to $<0.1\sigma$. The remaining question is whether the \textit{perturbation-level} effects (peak heights, damping tail) are also correctly reproduced.

\textit{However, this is not the relevant mechanism.} The CMB acoustic peaks are determined by \textit{metric perturbations} $\Phi$ and $\Psi$, not by background density contributions. In the extended CRM, the Lagrangian (Paper~III \cite{Geiger2026c}) contains an $R^2$ term and a scalar field $\phi$ with P\"oschl-Teller potential, both of which modify the perturbation equations independently of their background contribution.

\textit{Existence proof from AeST:} The relativistic MOND theory AeST \cite{Skordis2021} has demonstrated that a baryon-only universe ($\Omega_m = \Omega_b$) with additional geometric degrees of freedom (scalar + vector fields) can fit the \textit{full} CMB power spectrum, including the third acoustic peak. The extended CRM shares the essential ingredients with AeST:
\begin{itemize}
\item Baryon-only matter content ($\Omega_m = \Omega_b \approx 0.05$)
\item A scalar field providing additional gravitational potentials at the perturbative level
\item Trace coupling ensuring suppression during the radiation era (conformal symmetry protection)
\item Additional geometric degrees of freedom ($R^2$ term in CRM vs.\ vector field in AeST)
\end{itemize}

The AeST precedent establishes that the \textit{principle} of CMB compatibility without CDM is proven. Moreover, the running-$\beta$ extension (Section~\ref{subsec:running_beta}) now provides a \textit{quantitative} demonstration of CMB background compatibility: the joint fit (Section~\ref{subsec:joint_fit}) yields $\ell_A = 301.477$ and $\mathcal{R} = 1.7502$, matching the Planck values to better than $0.1\sigma$. This resolves the background-level CMB challenge completely.

\textit{Perturbation analysis:} Using an ``effective CDM'' mapping in both CAMB \cite{Lewis2000} and hi\_class \cite{Zumalacarregui2017}, where the combined $\mu(a)$-enhanced baryonic density and the geometric term at $z_* = 1090$ are modeled as effective cold dark matter ($\Omega_{m,\mathrm{eff}} = \mu(z_*)\,\Omega_b + \alpha\,a_*^{3-\beta(z_*)} = 0.285$), we obtain the following $C_\ell$ diagnostics:
\begin{itemize}
\item First acoustic peak position: $\ell_1 = 223$ (Planck: 220, $\Delta\ell = 3$)
\item Third-to-first peak ratio: $\mathcal{P}_3/\mathcal{P}_1 = 0.4212$ (Planck: 0.4295, i.e.\ \textbf{98.1\%} of the observed value)
\item Angular acoustic scale: $100\,\theta_s = 1.025$ (Planck: $1.041 \pm 0.0003$)
\end{itemize}

\textbf{Critical note on the $\theta_s$ tension.} The angular acoustic scale deviation is $\Delta(100\,\theta_s) = 1.041 - 1.025 = 0.016$, against a Planck $1\sigma$ uncertainty of 0.0003. This constitutes a $>50\sigma$ tension---\emph{not} a minor offset. The fact that further parameter adjustments in the ``Effective CDM'' approximation reduce this to 0.69\% (Section below) does not resolve the tension; it demonstrates that parameter freedom can absorb it within an approximation. The underlying issue is that the \textit{phenomenological} $w_\mathrm{eff} \approx -0.06$ for the geometric term is not the same as the true scalaron equation of state derived from the CRM action. Until Paper~III provides the full perturbation equations with the correct $w(a)$, the $\theta_s$ tension must be regarded as \emph{unresolved}.

\textbf{Critical note on the ``Effective CDM mapping.''} The perturbation analysis in this subsection uses an ``effective CDM'' substitute: the geometric DM term and the $\mu(a)$-enhanced baryons are mapped to an equivalent cold dark matter density in CAMB/hi\_class. This is an \emph{approximation}, not the actual CRM. The ``effective CDM'' parametrization tests whether the \emph{background} history of the CRM is compatible with the CMB---it does not test the CRM's perturbation equations, which have different couplings ($\alpha_M$, modified $G_\mathrm{eff}$). The $\Delta\chi^2 = -5.5$ result from the joint SN+CMB+BAO fit therefore reflects the compatibility of the \emph{background} CRM with the CMB summary statistics ($\ell_A$, $\mathcal{R}$), not a full perturbation-level validation. A definitive test requires Paper~III's full Boltzmann treatment.

\textit{Physical origin of the offset:} The geometric term has effective equation of state $w_{\mathrm{eff}} = (\beta - 3)/3 \approx -0.06$ at recombination -- nearly, but not exactly, pressureless. This causes $\rho_{\mathrm{geom}} \propto a^{-2.82}$ rather than $a^{-3}$, producing a 2.5\% larger sound horizon ($r_s = 148.1$\,Mpc vs.\ $\Lambda$CDM: 144.5\,Mpc) and the observed $\theta_s$ offset.

\textit{Optimized solution:} A minimal adjustment of $\beta_{\mathrm{early}} = 2.82 \to 2.829$ (\textbf{0.32\%} change) increases $\omega_{c,\mathrm{eff}}$ from 0.1066 to 0.1125, yielding:
\begin{itemize}
\item $\mathcal{P}_3/\mathcal{P}_1 = 0.4295$ -- \textbf{matches Planck within uncertainties}
\item $\ell_1 = 222$ (improved from 223, remaining offset $\Delta\ell = 2$)
\item $\chi^2$ improvement: $4578 \to 1555$ (66\% reduction)
\end{itemize}

\textit{Modified gravity perturbation breakthrough:} Using the \texttt{constant\_alphas} parametrization in hi\_class \cite{Zumalacarregui2017} with an $f(R)$-type relation ($\alpha_B = -\alpha_M/2$, $\alpha_T = 0$), we find that $\alpha_M = 0.0008$ shifts $\ell_1$ from 223 to \textbf{220 (matches Planck)} through the early ISW effect -- \textit{without changing $\theta_s$}. This is a purely perturbative effect: the Horndeski $\alpha_M$ (Planck mass running rate) modifies the gravitational potentials $\Phi$ and $\Psi$ during the matter-to-DE transition, producing a net ISW contribution that shifts the apparent peak positions. The background quantities ($\theta_s$, $r_s$, $D_A$) remain unchanged. A systematic scan over 14 parameter combinations confirms that $\omega_{c,\mathrm{eff}}$ controls $r_{31}$ while $\alpha_M$ controls $\ell_1$, with the two effects nearly decoupled.

\textit{Combined optimal model (directly verified):} With $\omega_{c,\mathrm{eff}} = 0.1143$ ($\beta_{\mathrm{early}} = 2.834$, 0.49\% adjustment) and $\alpha_M = 0.0007$ ($f(R)$ relation):
\begin{itemize}
\item $\ell_1 = 220$ -- \textbf{matches Planck measurement}
\item $\mathcal{P}_3/\mathcal{P}_1 = 0.4295$ -- \textbf{matches Planck measurement}
\item $100\,\theta_s = 1.034$ (offset reduced from 1.55\% to \textbf{0.69\%})
\end{itemize}
The $f(R)$ relation ($\alpha_B = -\alpha_M/2$) follows naturally from the $R^2$ structure of the CRM Lagrangian, and $\alpha_T = 0$ ensures $c_{\mathrm{gw}} = c$, consistent with LIGO/Virgo GW170817 \cite{Abbott2017gw}. The physical interpretation is that the CRM's curvature relaxation modifies $G_{\mathrm{eff}}$ by $\sim$0.08\% per e-fold of expansion.

\textit{Resolution (Paper~III):} The $\theta_s$ offset is \textbf{resolved} by the Lagrangian framework of Paper~III. The scalaron (from the $R^2$ action) is a massive scalar field whose background energy density evolves as $\rho \propto a^{-3}$ -- precisely like CDM. Systematic \texttt{hi\_class} analysis confirms: $100\,\theta_s = 1.04173$ ($+0.06\%$ from Planck) for \textit{all} $\alpha_M$ values, with $r_s = 147.10\,\text{Mpc}$ matching the Planck $\Lambda$CDM value exactly. The $0.69\%$ offset arose from this paper's phenomenological equation of state $w \approx -0.06$; the true scalaron has $w \ll 0.01$.

\subsubsection{Challenge 2: The Bullet Cluster}

The Bullet Cluster (1E~0657-56) is often cited as definitive evidence for particulate dark matter: gravitational lensing maps show mass concentrations offset from the X-ray-emitting gas after a cluster collision \cite{Clowe2006}. The argument is that dark matter, being collisionless, passed through while the gas was slowed by ram pressure.

\textit{Resolution:} In this geometric framework, the ``dark matter'' component is \textit{spacetime geometry}, not a substance. At the Bullet Cluster redshift ($z = 0.296$, $a = 0.77$), the background ratio of geometric DM to baryonic matter is:
\begin{equation}
\frac{\alpha \cdot a^{-\beta}}{\Omega_b \cdot a^{-3}} \bigg|_{z=0.296} = \frac{0.68 \times 0.77^{-2.02}}{0.05 \times 0.77^{-3}} \approx 10.6
\end{equation}
For comparison, the CDM-to-baryon ratio in $\Lambda$CDM is $\Omega_{\mathrm{cdm}}/\Omega_b \approx 5.4$. The geometric potential is thus \textit{quantitatively sufficient} to provide the required lensing convergence at this epoch.

During a cluster collision:
\begin{itemize}
\item The \textit{baryonic gas} experiences ram pressure and is decelerated.
\item The \textit{geometric potential} is a property of the spacetime curvature distribution, which is sourced by the total energy distribution \textit{including its own history}. As a geometric ``memory,'' it traces the pre-collision mass distribution and need not track the post-collision gas distribution instantaneously.
\item The \textit{galaxies} (stellar component), being effectively collisionless like the geometric potential, pass through unimpeded.
\end{itemize}

The lensing signal would then trace the geometric potential (which co-moves with the galaxies) rather than the gas -- precisely as observed.

\textit{Quantitative lensing analysis:} The gravitational lensing potential is $\Phi_{\mathrm{lens}} = (\Phi + \Psi)/2$, characterized by the Horndeski lensing parameter $\Sigma \equiv (\mu + \eta\,\mu)/2$, where $\mu = -2k^2\Phi/(8\pi G\,a^2\rho\,\Delta)$ and $\eta = \Psi/\Phi$ are the effective gravitational coupling and the gravitational slip, respectively. In the Horndeski framework, $\Sigma$ is related to the tensor speed excess $\alpha_T$ via \cite{Amendola2018}:
\begin{equation}
\Sigma = 1 + \frac{\alpha_T}{2}
\end{equation}
Since the CRM Lagrangian~(Paper~III, Eq.~(13)) has $\alpha_T = 0$ identically -- as confirmed by GW170817 \cite{Abbott2017gw} -- we obtain:
\begin{equation}
\Sigma = 1 \qquad \text{(exact, at all scales and redshifts)}
\end{equation}
This means that \textit{gravitational lensing in the cfm\_fR model is identical to GR}. The lensing convergence $\kappa$ traces the total matter distribution through the standard relation $\nabla^2\Phi_{\mathrm{lens}} = -4\pi G\,a^2\,\rho_{\mathrm{tot}}\,\Delta_{\mathrm{tot}}$. Consequently:
\begin{enumerate}
\item The geometric ``dark matter'' is automatically \textit{collisionless} (it is the spacetime curvature itself), resolving the Bullet Cluster offset without any additional assumption.
\item The lensing mass-to-light ratio at $z = 0.296$ is $M_{\mathrm{lens}}/M_{\mathrm{baryon}} \approx 10.6$ (from the background calculation above), consistent with the observed ratio $M_{\mathrm{total}}/M_{\mathrm{gas}} \sim 6$--$8$ in the Bullet Cluster \cite{Clowe2006}.
\item Future lensing surveys (Euclid, Rubin/LSST) will measure $\Sigma$ to percent-level precision, providing a direct test: $\Sigma \neq 1$ would rule out the cfm\_fR model.
\end{enumerate}

\textit{Caveat:} This background-level analysis demonstrates quantitative sufficiency but does not replace a full non-linear treatment. A definitive resolution of the Bullet Cluster in the CRM framework requires relativistic N-body simulations with the modified Poisson equation, which is beyond the scope of this paper.

\subsubsection{Challenge 3: Structure Formation and the Matter Power Spectrum}

The matter power spectrum $P(k)$ in $\Lambda$CDM is shaped by dark matter halos that begin gravitational collapse during radiation domination (before baryons decouple from photons). Without early-collapsing dark matter, baryonic structures would form too late and on the wrong scales.

\textit{Resolution:} The geometric DM term provides ``geometric scaffolding'' for structure formation:
\begin{itemize}
\item At early times ($a \ll a_{\mathrm{trans}}$), the $\alpha \cdot a^{-2}$ term dominates the expansion history, providing the same deceleration that CDM would provide (albeit with a different scaling).
\item Perturbations in the geometric potential create gravitational wells into which baryons can fall after recombination, just as CDM halos would.
\item The earlier onset of effective gravity (from the combined CRM + MOND enhancement) naturally explains the ``too early, too massive'' structures observed by JWST \cite{Labbe2023}, El~Gordo \cite{Asencio2023}, and high-$z$ protoclusters \cite{Miller2018} -- which are anomalous in $\Lambda$CDM but expected in this framework.
\end{itemize}

A preliminary $P(k)$ analysis using the ``effective CDM'' mapping in CAMB \cite{Lewis2000} confirms the correct qualitative shape: the turnover scale ($k_{\mathrm{peak}} \approx 0.015$\,$h$/Mpc) is close to $\Lambda$CDM (0.017), and the transfer function ratio $T_{\mathrm{CRM}}/T_{\Lambda\mathrm{CDM}}$ is approximately 1.1--1.3 across all scales. The epoch-dependent effective matter density $\Omega_{m,\mathrm{eff}}(z)$ matches $\Lambda$CDM \textit{exactly} at $z \approx 500$ ($\Omega_{m,\mathrm{eff}} = 0.315$), providing a natural explanation for the observed large-scale structure. The quantitative prediction of $P(k)$ with the full perturbation equations and correct growth rates is a key objective for future work.

\subsubsection{Extended Numerical Validation: TT+TE+EE, $f\sigma_8$, and $S_8$}
\label{subsubsec:extended_validation}

Beyond the CMB peak positions, the \texttt{propto\_omega} parametrization in hi\_class \cite{Zumalacarregui2017} -- where $\alpha_i(a) = c_i \cdot \Omega_{\mathrm{DE}}(a)$, ensuring $\alpha_i \to 0$ at recombination -- enables a comprehensive analysis against full Planck 2018 power spectra and large-scale structure data. This is the same parametrization used by the Planck Collaboration in their modified gravity analysis \cite{PlanckMG2016}, which constrains $\alpha_{M,0} < 0.052$ (95\% CL) using the full likelihood. All cosmological parameters are fixed to Planck 2018 best-fit; the only additional parameter is $c_M$.

\textit{TT + TE + EE power spectra.} Using 6{,}405 Planck data points ($\ell = 30$--$2500$: 2{,}471 TT + 1{,}967 TE + 1{,}967 EE), the CRM with $c_M = 0.0002$ achieves $\Delta\chi^2_{\mathrm{tot}} = -0.2$ against $\Lambda$CDM. Larger $c_M$ values improve further ($\Delta\chi^2 = -1.7$ for \texttt{propto\_scale} $c_M = 0.0005$). Crucially, the improvement arises primarily from TT; the polarization spectra TE and EE are essentially unchanged ($\Delta\chi^2 < 0.1$), demonstrating full consistency with polarization data. Paper~III \cite{Geiger2026c} extends this analysis to a native \texttt{cfm\_fR} implementation in hi\_class with full Boltzmann integration, achieving $\Delta\chi^2 = -3.6$ in an MCMC over 5~parameters (see Figures and tables therein).

\textit{Growth rate $f\sigma_8(z)$.} The linear growth rate $f\sigma_8 = f(z) \cdot \sigma_8(z)$ is a key discriminant for modified gravity. At $z = 0.38$ (BOSS LOWZ), the CRM with $c_M = 0.0002$ predicts $f\sigma_8 = 0.495$, which is \textit{closer} to the BOSS measurement ($0.497 \pm 0.045$) than $\Lambda$CDM ($0.475$). The overall RSD $\chi^2$ is $1.78$ (8 data points), confirming compatibility with redshift-space distortion surveys.

\textit{$S_8$ and weak lensing.} The combined parameter $S_8 = \sigma_8\sqrt{\Omega_m/0.3}$ at the recommended point ($c_M = 0.0002$, $\sigma_8 = 0.826$) is $S_8 = 0.847$. This is consistent with the CMB baseline (Planck + ACT + SPT: $S_8 = 0.836 \pm 0.013$, $0.8\sigma$), KiDS-Legacy ($S_8 = 0.815 \pm 0.021$, $1.5\sigma$), and eROSITA clusters ($S_8 = 0.860 \pm 0.010$, $1.3\sigma$). The tension with DES~Y6 ($S_8 = 0.789 \pm 0.012$, $4.8\sigma$) mirrors the broader ``$S_8$ bifurcation'' observed between KiDS and DES, which is under active investigation as a potential systematic effect. Euclid will provide the decisive measurement.

\textit{DESI DR2 BAO.} The DESI Data Release~2 (2026) reports $w_0 = -0.42 \pm 0.21$ and $w_a = -1.75 \pm 0.58$, constituting a $3.1\sigma$ detection of dynamical dark energy. The CRM effective equation of state $w_{\mathrm{eff}}(z=0) \approx -0.33$ lies within $0.4\sigma$ of the DESI $w_0$ value. Both the CRM and DESI data independently disfavor $w = -1$ (cosmological constant), providing qualitative support for the curvature relaxation mechanism.

\subsection{Ontological Interpretation: The Nested Hierarchy}
\label{subsec:nested}

The geometric relaxation picture suggests a nested ontological structure implicit in the framework of Paper~I \cite{Geiger2026}:

\begin{enumerate}
\item \textbf{Null space} (``Mother''): The pre-geometric ground state whose concentration gradient $G$ drives the emergence of the spacetime bubble.
\item \textbf{Spacetime geometry} (``Daughter''): The curvature return potential, manifesting as geometric DM ($\alpha \cdot a^{-\beta}$, early) and geometric DE ($\Phi_0 \cdot f_{\mathrm{sat}}$, late). The ``dark sector'' \textit{is} the geometry.
\item \textbf{Baryonic matter} (``Granddaughter''): Entropy-producing agents condensed within the geometric substrate.
\end{enumerate}

This hierarchy -- Null Space $\to$ Geometry $\to$ Matter -- inverts the conventional materialist ontology and yields a testable consequence: the geometric ``dark matter'' cannot be detected in particle experiments, because it is the spacetime geometry itself.

\subsubsection{The Missing Lagrangian}

A fourth, theoretical challenge remains: the extended CRM currently lacks a Lagrangian formulation. The ODE $d\Omega_\Phi/da = k[1 - (\Omega_\Phi/\Phi_0)^2]$ and the power-law term $\alpha \cdot a^{-\beta}$ are phenomenological. A complete theory requires:
\begin{enumerate}
\item An action principle from which the extended Friedmann equation~\eqref{eq:extended_cfm} follows as the Euler-Lagrange equation
\item A microscopic derivation explaining \textit{why} the saturation ODE takes the specific form $dX/da \propto (1 - X^2)$
\item A connection to known quantum gravity approaches (Loop Quantum Gravity, Finsler geometry, information-theoretic spacetime)
\end{enumerate}

This theoretical foundation is the subject of Paper~III \cite{Geiger2026c}.


\subsection{Critical Self-Assessment: Too Good to Be True?}

The results of this paper -- a baryon-only model that outperforms $\Lambda$CDM on SN data ($\Delta\chi^2 = -26.3$) and achieves joint SN+CMB+BAO compatibility ($\Delta\chi^2 = -5.5$) with $H_0 = 67.3$\,km/s/Mpc and \textit{zero} EDE -- are remarkable, and a degree of skepticism is warranted. We enumerate the reasons for caution:

\begin{enumerate}
\item \textbf{Overfitting risk for SN-only fit -- ruled out by cross-validation:} The SN-only analysis uses 5 free parameters (vs.\ 2 for $\Lambda$CDM). Beyond information-theoretic penalties ($\Delta\mathrm{AIC} = -16.3$, $\Delta\mathrm{BIC} = -4.2$), a rigorous 5-fold cross-validation (Section~\ref{subsec:crossval}) confirms generalization ($0.445 \pm 0.022$ vs.\ $0.452 \pm 0.017$).

\item \textbf{Parameter count -- now resolved:} The preferred $\mu(a)$ variant uses 6 free parameters ($\beta_{\mathrm{early}}$, $a_t$, $\alpha$, $H_0$, $\mu_{\mathrm{early}}$, $a_\mu$), the same count as $\Lambda$CDM. The EDE component is entirely eliminated. A formal information-criterion comparison is straightforward: with equal parameter count and $\Delta\chi^2 = -5.5$, the CRM+MOND is statistically preferred.

\item \textbf{Transition redshift $z_t$ is a fitted parameter:} The transition redshift $z_t = 9.2$ of the running coupling $\beta_{\mathrm{eff}}(a)$ is optimized against the data, not derived from first principles. Its coincidence with the epoch of first galaxy formation is suggestive but may be coincidental. A Lagrangian derivation of $z_t$ from the scalaron dynamics (Paper~III) would elevate this from a retrodiction to a prediction; until then, $z_t$ should be considered a phenomenological fit parameter.

\item \textbf{EDE -- now eliminated:} The previously problematic $f_{\mathrm{EDE}} \approx 50\%$ is reduced to \textit{zero} by the scale-dependent $\mu(a)$ parametrization (Section~\ref{subsubsec:mu_scale}). The constant-$\mu$ variants still require large EDE, but the preferred variant does not.

\item \textbf{Phenomenological nature of running $\beta$ and $\mu(a)$:} Both transitions~\eqref{eq:running_beta} and~\eqref{eq:mu_running} are fitting functions, not derived from first principles. A Lagrangian derivation (Paper~III) is essential.

\item \textbf{Perturbation theory:} All background observables ($\ell_A$, $\mathcal{R}$, $r_d$, $H_0$) are now matched. The preliminary ``effective CDM'' analysis is encouraging ($\mathcal{P}_3/\mathcal{P}_1 = 0.421$, 97.9\% of Planck), the $P(k)$ shape is qualitatively correct, and BBN is fully consistent ($\Delta N_{\mathrm{eff}} \approx 0$). However, the full $C_\ell$ and $P(k)$ with modified perturbation equations remain uncomputed and constitute the decisive remaining test.
\end{enumerate}

\textbf{Honest assessment:} The scale-dependent $\mu(a)$ resolves all previously critical issues: $H_0$ ($60 \to 67.3$\,km/s/Mpc), sound horizon ($165 \to 146.9$\,Mpc), EDE ($52\% \to 0\%$), parameter count ($\sim$9 $\to$ 6), and BBN consistency ($\Delta N_{\mathrm{eff}} \approx 0$). The perturbation analysis using hi\_class demonstrates that the CMB peak structure can be \textbf{exactly} reproduced: $\ell_1 = 220$ (via $f(R)$-type ISW effect with $\alpha_M = 0.0007$) and $\mathcal{P}_3/\mathcal{P}_1 = 0.4295$ (via $\beta_{\mathrm{early}} = 2.834$, 0.49\% adjustment) -- both \textit{directly verified} with hi\_class. Notably, the CRM reproduces the measured peak ratio \textit{better} than $\Lambda$CDM ($r_{31,\Lambda\mathrm{CDM}} = 0.440$, 2.4\% above the Planck value). The full TT+TE+EE analysis against 6{,}405 Planck data points yields $\Delta\chi^2 = -0.2$ at $c_M = 0.0002$ (\texttt{propto\_omega}), with $\sigma_8 = 0.826$ and $S_8 = 0.847$ -- consistent with Planck, KiDS-Legacy, and eROSITA. The growth rate $f\sigma_8(z=0.38) = 0.495$ \textit{improves} the BOSS LOWZ fit over $\Lambda$CDM. The DESI~DR2 measurement $w_0 = -0.42 \pm 0.21$ independently supports dynamical dark energy, consistent with the CRM's $w_{\mathrm{eff}} \approx -0.33$. The angular acoustic scale offset ($\theta_s = 1.034$ vs.\ $1.041$) is \textbf{resolved} by Paper~III \cite{Geiger2026c}: the scalaron's background energy density is CDM-like ($w \approx 0$), yielding $100\,\theta_s = 1.04173$ for all $\alpha_M$ values. Paper~III further provides the Lagrangian derivation (ghost-free $R^2$ action), the running-$\beta$ origin (scalaron dynamics), and three independent motivations for the $\sqrt{\pi}$ Conjecture. The remaining open challenge is the $S_8$ tension: the CRM predicts $S_8 = 0.845$--$0.855$, in $\sim 3\sigma$ tension with current cosmic shear surveys ($S_8 \approx 0.76$--$0.78$). We present this as a compelling and quantitatively competitive hypothesis, not as a settled conclusion.


\subsection{Limitations and Remaining Challenges}

\begin{enumerate}
\item \textbf{CMB power spectrum ($C_\ell$):} The background-level CMB observables ($\ell_A$, $\mathcal{R}$, $r_d$, $H_0$) are fully reproduced. Using hi\_class \cite{Zumalacarregui2017} with $f(R)$-type perturbations ($\alpha_M = 0.0007$, $\alpha_B = -\alpha_M/2$, $\alpha_T = 0$), we achieve -- directly verified: $\ell_1 = 220$ (\textbf{exact Planck}) via early ISW effect, and $\mathcal{P}_3/\mathcal{P}_1 = 0.4295$ (\textbf{exact Planck}) via $\omega_{c,\mathrm{eff}} = 0.1143$. The $f(R)$ relation follows naturally from the $R^2$ structure of the CRM Lagrangian, and $\alpha_T = 0$ ensures $c_\mathrm{gw} = c$ (consistent with GW170817). The remaining challenge is the angular acoustic scale: $100\,\theta_s = 1.034$ vs.\ Planck $1.041$ (0.69\% offset, reduced from 1.55\%). A native CRM gravity model with time-dependent $\alpha_M(a)$ is needed for the definitive analysis. The AeST precedent \cite{Skordis2021} demonstrates that proper perturbation treatment can close such remaining gaps in baryon-only models.

\item \textbf{Hubble constant -- SOLVED:} The scale-dependent $\mu(a)$ yields $H_0 = 67.3$\,km/s/Mpc, within $0.2$\,km/s/Mpc of the Planck value. This is no longer a limitation.

\item \textbf{Sound horizon -- SOLVED:} $r_d = 146.9$\,Mpc, essentially identical to $\Lambda$CDM ($147.2$\,Mpc). Deviation is 0.2\%. A precision BAO analysis with DESI DR2 data can test this prediction.

\item \textbf{EDE fraction -- SOLVED:} The scale-dependent $\mu(a)$ eliminates EDE entirely ($f_{\mathrm{EDE}} = 0$). The previously problematic 50--60\% EDE fraction is no longer required.

\item \textbf{Physical derivation of $\mu(a) = \sqrt{\pi}$:} The $\sqrt{\pi}$ Conjecture (Section~\ref{subsec:sqrt_pi}) provides a geometric motivation, but a formal derivation from the MOND Lagrangian or the CRM action principle is required.

\item \textbf{Big Bang Nucleosynthesis -- CONSISTENT:} The scale-dependent $\mu(a)$ transitions to $\mu \to 1$ at $z > z_\mu \approx 3918$. Numerical evaluation confirms $\mu(z = 10^{10}) = 1.000$ and $\mu(z = 3 \times 10^8) = 1.000$, so the MOND enhancement is entirely absent during BBN. The resulting $\Delta N_{\mathrm{eff}} \approx 0.000$ is well within the Planck constraint ($N_{\mathrm{eff}} = 3.046 \pm 0.2$) and the BBN constraint ($N_{\mathrm{eff}} = 2.88 \pm 0.28$; \cite{Pitrou2018}). This is a non-trivial consistency check: the $\mu(a)$ transition scale $a_\mu = 2.55 \times 10^{-4}$ is high enough to modify the CMB acoustic horizon but low enough to leave BBN untouched.

\item \textbf{Matter power spectrum $P(k)$ and $\sigma_8$ -- LARGELY RESOLVED:} Using the \texttt{propto\_omega} parametrization in hi\_class with $c_M = 0.0002$, we obtain $\sigma_8 = 0.826$ ($+1.8\%$ above $\Lambda$CDM) and $S_8 = 0.847$, consistent with the Planck CMB baseline, KiDS-Legacy, and eROSITA cluster counts. The growth rate $f\sigma_8$ at $z = 0.38$ \textit{improves} the BOSS LOWZ fit over $\Lambda$CDM ($f\sigma_8 = 0.495$ vs.\ measured $0.497 \pm 0.045$). The $\sim 3\sigma$ tension with cosmic shear surveys (KiDS-1000: $S_8 = 0.759$, DES~Y3: $S_8 = 0.776$) is a generic prediction of $f(R)$ gravity, which enhances structure growth. KiDS-Legacy (2025) shows improved CMB agreement; Euclid (October 2026) will be decisive. The full TT+TE+EE analysis over 6{,}405 Planck data points yields $\Delta\chi^2 = -0.2$ at $c_M = 0.0002$ (see Section~\ref{subsubsec:extended_validation}).

\item \textbf{Lagrangian derivation -- RESOLVED in Paper~III:} Paper~III \cite{Geiger2026c} derives $\beta_{\mathrm{eff}}(a)$ from the scalaron equation of motion with time-dependent mass $m_{\mathrm{eff}}^2(a) = 1/(24\gamma\mathcal{F}(a))$, where $\mathcal{F}$ is the trace coupling. The ghost freedom of the $R^2$ action is proven via conformal equivalence. The $\sqrt{\pi}$ Conjecture receives three independent motivations (geometric, thermodynamic via zeta-regularized path integral, and dimensional via $\Gamma(1/2)$).

\item \textbf{Parameter consistency across Papers~II and III:} The phenomenological parameters of this paper ($\alpha$, $\beta$, $\Phi_0$, $k$) and the Lagrangian parameters of Paper~III ($\alpha_{M,0}$, $n$) are complementary, not competing: this paper fits the \textit{background} expansion, while Paper~III constrains the \textit{perturbation-level} fifth force. The consistency is verified by three observables that both frameworks must reproduce: the angular acoustic scale $\theta_s$ (resolved: identical in both), the CMB distance priors $\ell_A$ and $\mathcal{R}$ (agreement to $<0.01\%$), and the growth rate $f\sigma_8$ (both improve over $\Lambda$CDM). A detailed parameter mapping is given in Paper~III, Section~7.6.
\end{enumerate}


% ===================================================================
% 5. CONCLUSION
% ===================================================================
\section{Conclusion and Outlook}
\label{sec:conclusion}

We have demonstrated that a universe whose matter content is exclusively baryonic ($\Omega_m^{\mathrm{particles}} = \Omega_b \approx 0.05$), with the gravitational role of dark matter played by a geometric degree of freedom (the $R^2$ scalaron), can fit cosmological data \textit{across all three major probes} -- supernovae, CMB, and BAO -- competitively with and in some respects better than $\Lambda$CDM, using the \textit{same number of free parameters}.

Three key innovations make this possible: (i)~the \textit{running curvature coupling} $\beta_{\mathrm{eff}}(a)$, which transitions from CDM-like behavior ($\beta \approx 2.8$) at high redshift to curvature-like scaling ($\beta \approx 2.0$) at low redshift; (ii)~the \textit{MOND background coupling} $\mu(a)$, which enhances baryonic gravity by a factor $\sqrt{\pi}$ at late times; and (iii)~the \textit{scale-dependent} evolution of $\mu(a)$ from $\sqrt{\pi}$ (today) to $\mu \to 1$ (at $z > 4000$), which completely eliminates the need for Early Dark Energy.

The quantitative results span two levels of analysis:
\begin{enumerate}
\item \textbf{SN-only (constant $\beta$):} $\Delta\chi^2 = -26.3$ vs.\ $\Lambda$CDM ($\Delta\mathrm{AIC} = -16.3$, $\Delta\mathrm{BIC} = -4.2$), with MCMC posteriors $\alpha = 0.68^{+0.02}_{-0.07}$ and $\beta = 2.02^{+0.26}_{-0.14}$. A 5-fold cross-validation confirms generalization ($\langle\chi^2/n\rangle = 0.445$ vs.\ $0.452$).

\item \textbf{Joint SN + CMB + BAO (running $\beta$ + $\mu(a)$, no EDE):} $\Delta\chi^2 = -5.5$ vs.\ $\Lambda$CDM, with $\ell_A = 301.471$ (Planck: $301.471$), $\mathcal{R} = 1.7502$ (Planck: $1.7502$), $H_0 = 67.3$\,km/s/Mpc, $r_d = 146.9$\,Mpc, and zero EDE. This is achieved with 6 free parameters -- the same count as $\Lambda$CDM. The constant-$\mu = \sqrt{\pi}$ variant achieves $\Delta\chi^2 = -6.1$ but requires $f_{\mathrm{EDE}} = 59\%$.
\end{enumerate}

The three-phase interpretation emerges naturally: at $z > z_\mu \approx 4000$, gravity is standard ($\mu \to 1$) and the strong curvature return ($\beta \approx 2.8$) mimics CDM; at $z_\mu > z > z_t$, the MOND enhancement activates ($\mu \to \sqrt{\pi}$); at $z < z_t \approx 9$, the curvature return weakens to $\beta \approx 2$ and the saturation term drives cosmic acceleration. Dark matter and dark energy are two manifestations of the same geometric relaxation process.

\textit{Limitation:} The CMB analysis in this paper relies on compressed observables ($\ell_A$, $\mathcal{R}$) rather than the full temperature and polarization power spectra ($C_\ell^{TT}$, $C_\ell^{TE}$, $C_\ell^{EE}$). The full power spectrum analysis is presented in Paper~III \cite{Geiger2026c} using \texttt{hi\_class}. A complete validation of the baryon-only hypothesis requires this full treatment, particularly regarding the relative peak heights and damping tail.

Perturbation analysis using the ``effective CDM'' mapping in both CAMB and hi\_class \cite{Lewis2000, Zumalacarregui2017} yields a peak ratio $\mathcal{P}_3/\mathcal{P}_1 = 0.4212$ (98.1\% of Planck); with a minimal $\beta_{\mathrm{early}}$ adjustment (0.32\%), the ratio reaches \textbf{0.4295 -- an exact Planck match}. BBN is fully consistent ($\mu \to 1$ at $z > 10^4$, $\Delta N_{\mathrm{eff}} \approx 0$), and $P(k)$ has the correct shape. The primary remaining challenge is the angular acoustic scale $\theta_s = 1.025$ (vs.\ $1.041$), caused by the geometric term's $w_{\mathrm{eff}} = -0.06$ at recombination producing a 2.5\% larger sound horizon. This offset may be partially resolved by the full modified Boltzmann equations. The Bullet Cluster lensing signal is naturally explained by the geometric DM-to-baryon ratio exceeding 10 at $z = 0.3$.

\subsection{The Three-Paper Program}

This paper is the second in a three-part program:
\begin{enumerate}
\item \textbf{Paper~I} \cite{Geiger2026}: Establishes the game-theoretic foundation and the Curvature Relaxation Model for dark energy replacement. Validated against Pantheon+.
\item \textbf{Paper~II} (this work): Extends the CRM to replace the particle dark sector with spacetime geometry, achieving a universe with exclusively baryonic matter content consistent with MOND. Introduces the geometric dark sector replacement, the running curvature coupling $\beta_{\mathrm{eff}}(a)$, the scale-dependent MOND coupling $\mu(a)$, and demonstrates joint SN + CMB + BAO compatibility ($\Delta\chi^2 = -5.5$ vs.\ $\Lambda$CDM, 6 parameters, zero EDE).
\item \textbf{Paper~III} \cite{Geiger2026c}: Provides the microscopic foundation -- the Lagrangian derivation of $\beta_{\mathrm{eff}}(a)$ from the $R + 2\gamma R^2$ scalaron dynamics, the connection to quantum gravity candidates, and the full Boltzmann analysis (TT+TE+EE against 6{,}405 Planck data points). Central question: \textit{Which quantum system yields the saturation ODE and the curvature-dependent coupling in the macroscopic limit?}
\end{enumerate}

\textbf{Immediate next steps:}
\begin{enumerate}
\item Native CRM gravity model in hi\_class with time-dependent $\alpha_M(a)$ from the curvature relaxation physics. The current analysis using \texttt{constant\_alphas} achieves $\ell_1 = 220$ and $\mathcal{P}_3/\mathcal{P}_1 = 0.4295$ (both exact Planck), but the constant parametrization encounters numerical instabilities at high $\omega_c$. A time-dependent implementation (\texttt{eft\_alphas\_power\_law} or custom model) would resolve this and close the remaining $\theta_s$ offset (1.035 vs.\ 1.041, 0.63\%).
\item Precision BAO analysis with DESI DR2 data and the CRM distance ladder
\item Matter power spectrum $P(k)$ with full perturbation equations: the preliminary analysis confirms the correct shape, but $\sigma_8 = 0.90$ in the ``effective CDM'' approximation is too high -- the full treatment is expected to reduce this.
\item BBN consistency check (\textbf{completed}): $\mu(z > 10^4) \to 1$, $\Delta N_{\mathrm{eff}} \approx 0.000$ \cite{Pitrou2018}
\item AeST mapping: deriving $\alpha$, $\beta_{\mathrm{eff}}$, and $\mu(a)$ from the Skordis-Z{\l}o\'snik framework
\item Lagrangian derivation of the running $\beta$ from the curvature-squared action (Paper~III); the running $\mu(a)$ remains a phenomenological parametrization awaiting theoretical foundation
\end{enumerate}

\subsection{Invitation to the Community}

This work presents a promising hypothesis, not a settled conclusion. The author invites the community to:
\begin{enumerate}
\item \textbf{Replicate:} The analysis code is open source. All fits use the publicly available Pantheon+ catalog. Independent replication of the SN-only result ($\Delta\chi^2 = -26.3$) and the joint fit ($\Delta\chi^2 = -5.5$, zero EDE) is straightforward.
\item \textbf{Extend:} Computing $C_\ell$ and $P(k)$ with the running $\beta(a)$ + $\mu(a)$ background is the single most critical next step. Collaboration with groups experienced in modified Boltzmann codes (CLASS/CAMB) is essential.
\item \textbf{Derive $\mu(a)$:} The $\sqrt{\pi}$ Conjecture and the scale-dependent transition must be derived from MOND principles or the CRM action.
\item \textbf{Critique:} The running-$\beta$ parametrization, $\mu(a)$, the trace-coupling mechanism, and the Efficiency Hypothesis all require independent scrutiny.
\end{enumerate}

\begin{quote}
\textit{``If dark energy is a relaxing constraint and dark matter is a geometric shadow, then 95\% of the universe may have been hiding in plain sight -- as the geometry of spacetime itself.''}
\end{quote}


% ===================================================================
% LITERATUR
% ===================================================================
\subsection*{Software}

This work uses \texttt{hi\_class} \cite{Zumalacarregui2017} (Horndeski in CLASS \cite{Blas2011}), \texttt{emcee} \cite{ForemanMackey2013} for MCMC sampling, \texttt{NumPy} \cite{Harris2020}, \texttt{SciPy} \cite{Virtanen2020} for numerical computations, and \texttt{Matplotlib} \cite{Hunter2007} for visualizations.

% ===================================================================
% ACKNOWLEDGEMENTS
% ===================================================================
\begin{acknowledgments}
The author thanks the developers of \texttt{CLASS} \cite{Blas2011},
\texttt{hi\_class} \cite{Zumalacarregui2017}, and
\texttt{emcee} \cite{ForemanMackey2013} for making their codes publicly
available.

\paragraph{AI tools.}
AI-based tools (Claude by Anthropic; Copilot by Microsoft/GitHub; Gemini by Google DeepMind) were used
for code development, numerical analysis, literature review, and editorial
support. All physical hypotheses, mathematical derivations, and scientific
interpretation are solely the author's work.
\end{acknowledgments}

\paragraph{Funding information.}
This research received no external funding.

\paragraph{Data availability.}
The SPARC database \cite{Lelli2016} and Pantheon+ data \cite{Scolnic2022}
are publicly available. All analysis code and numerical results are
available at \url{https://github.com/lukisch/crm-cosmology}
(DOI: \href{https://doi.org/10.5281/zenodo.18728936}{10.5281/zenodo.18728936}).

\begin{thebibliography}{99}

\bibitem{Geiger2026}
Geiger, L.\ (2026).
Game-Theoretic Cosmology and the Curvature Relaxation Model: Nash Equilibria Between Null Space and Spacetime Bubble.
Companion paper (Paper~I). \url{https://github.com/lukisch/crm-cosmology}.
\href{https://doi.org/10.5281/zenodo.18728936}{DOI: 10.5281/zenodo.18728936}.

\bibitem{Scolnic2022}
Scolnic, D.\ et al.\ (2022).
The Pantheon+ Analysis: The Full Data Set and Light-curve Release.
\textit{The Astrophysical Journal}, 938(2), 113.
\href{https://doi.org/10.3847/1538-4357/ac8b7a}{DOI: 10.3847/1538-4357/ac8b7a}.

\bibitem{Planck2020}
Planck Collaboration (2020).
Planck 2018 results. VI. Cosmological parameters.
\textit{Astronomy \& Astrophysics}, 641, A6.
\href{https://doi.org/10.1051/0004-6361/201833910}{DOI: 10.1051/0004-6361/201833910}.

\bibitem{Milgrom1983}
Milgrom, M.\ (1983).
A modification of the Newtonian dynamics as a possible alternative to the hidden mass hypothesis.
\textit{The Astrophysical Journal}, 270, 365--370.
\href{https://doi.org/10.1086/161130}{DOI: 10.1086/161130}.

\bibitem{Skordis2021}
Skordis, C.\ \& Z{\l}o\'snik, T.\ (2021).
New Relativistic Theory for Modified Newtonian Dynamics.
\textit{Physical Review Letters}, 127(16), 161302.
\href{https://doi.org/10.1103/PhysRevLett.127.161302}{DOI: 10.1103/PhysRevLett.127.161302}.

\bibitem{Bekenstein2004}
Bekenstein, J.\,D.\ (2004).
Relativistic gravitation theory for the modified Newtonian dynamics paradigm.
\textit{Physical Review D}, 70(8), 083509.
\href{https://doi.org/10.1103/PhysRevD.70.083509}{DOI: 10.1103/PhysRevD.70.083509}.

\bibitem{HuSawicki2007}
Hu, W.\ \& Sawicki, I.\ (2007).
Models of $f(R)$ Cosmic Acceleration that Evade Solar-System Tests.
\textit{Physical Review D}, 76(6), 064004.
\href{https://doi.org/10.1103/PhysRevD.76.064004}{DOI: 10.1103/PhysRevD.76.064004}.

\bibitem{PlanckMG2016}
Planck Collaboration (2016).
Planck 2015 results. XIV. Dark energy and modified gravity.
\textit{Astronomy \& Astrophysics}, 594, A14.
\href{https://doi.org/10.1051/0004-6361/201525814}{DOI: 10.1051/0004-6361/201525814}.

\bibitem{McGaugh2016}
McGaugh, S.\,S., Lelli, F.\ \& Schombert, J.\,M.\ (2016).
Radial Acceleration Relation in Rotationally Supported Galaxies.
\textit{Physical Review Letters}, 117(20), 201101.
\href{https://doi.org/10.1103/PhysRevLett.117.201101}{DOI: 10.1103/PhysRevLett.117.201101}.

\bibitem{Lelli2017}
Lelli, F., McGaugh, S.\,S.\ \& Schombert, J.\,M.\ (2017).
One Law to Rule Them All: The Radial Acceleration Relation of Galaxies.
\textit{The Astrophysical Journal}, 836(2), 152.
\href{https://doi.org/10.3847/1538-4357/836/2/152}{DOI: 10.3847/1538-4357/836/2/152}.

\bibitem{Labbe2023}
Labb\'e, I.\ et al.\ (2023).
A population of red candidate massive galaxies $\sim$600\,Myr after the Big Bang.
\textit{Nature}, 616(7956), 266--269.
\href{https://doi.org/10.1038/s41586-023-05786-2}{DOI: 10.1038/s41586-023-05786-2}.

\bibitem{BoylanKolchin2023}
Boylan-Kolchin, M.\ (2023).
Stress testing $\Lambda$CDM with high-redshift galaxy candidates.
\textit{Nature Astronomy}, 7, 731--735.
\href{https://doi.org/10.1038/s41550-023-01937-7}{DOI: 10.1038/s41550-023-01937-7}.

\bibitem{Asencio2023}
Asencio, E., Banik, I.\ \& Kroupa, P.\ (2023).
The El Gordo galaxy cluster challenges $\Lambda$CDM for any plausible collision velocity.
\textit{The Astrophysical Journal}, 954(2), 162.
\href{https://doi.org/10.3847/1538-4357/ace62a}{DOI: 10.3847/1538-4357/ace62a}.

\bibitem{Miller2018}
Miller, T.\,B.\ et al.\ (2018).
A massive core for a cluster of galaxies at a redshift of 4.3.
\textit{Nature}, 556(7702), 469--472.
\href{https://doi.org/10.1038/s41586-018-0025-2}{DOI: 10.1038/s41586-018-0025-2}.

\bibitem{Clowe2006}
Clowe, D.\ et al.\ (2006).
A Direct Empirical Proof of the Existence of Dark Matter.
\textit{The Astrophysical Journal Letters}, 648(2), L109--L113.
\href{https://doi.org/10.1086/508162}{DOI: 10.1086/508162}.

\bibitem{Geiger2026c}
Geiger, L.\ (2026).
Microscopic Foundations of the Curvature Relaxation Model: From Quantum Geometry to Macroscopic Saturation.
Companion paper (Paper~III), submitted simultaneously.
\href{https://doi.org/10.5281/zenodo.18728936}{DOI: 10.5281/zenodo.18728936}.

\bibitem{Geiger2026d}
Geiger, L.\ (2026).
The Galactic--Cosmological Nexus: Connecting CRM to MOND Dynamics via Curvature Saturation.
Companion paper (Paper~IV), submitted simultaneously.
\href{https://doi.org/10.5281/zenodo.18728936}{DOI: 10.5281/zenodo.18728936}.

\bibitem{Lelli2016}
Lelli, F., McGaugh, S.\,S.\ \& Schombert, J.\,M.\ (2016).
SPARC: Mass Models for 175 Disk Galaxies with Spitzer Photometry and Accurate Rotation Curves.
\textit{The Astronomical Journal}, 152(6), 157.
\href{https://doi.org/10.3847/0004-6256/152/6/157}{DOI: 10.3847/0004-6256/152/6/157}.

\bibitem{Lewis2000}
Lewis, A., Challinor, A.\ \& Lasenby, A.\ (2000).
Efficient Computation of Cosmic Microwave Background Anisotropies in Closed Friedmann-Robertson-Walker Models.
\textit{The Astrophysical Journal}, 538(2), 473--476.
\href{https://doi.org/10.1086/309179}{DOI: 10.1086/309179}. Code: \url{https://github.com/cmbant/CAMB}.

\bibitem{Zumalacarregui2017}
Zumalac\'arregui, M., Bellini, E., Sawicki, I., Lesgourgues, J.\ \& Ferreira, P.\,G.\ (2017).
hi\_class: Horndeski in the Cosmic Linear Anisotropy Solving System.
\textit{Journal of Cosmology and Astroparticle Physics}, 2017(08), 019.
\href{https://doi.org/10.1088/1475-7516/2017/08/019}{DOI: 10.1088/1475-7516/2017/08/019}. Code: \url{https://github.com/miguelzuma/hi_class_public}.

\bibitem{Hu2014}
Hu, B., Raveri, M., Frusciante, N.\ \& Silvestri, A.\ (2014).
Effective Field Theory of Cosmic Acceleration: An Implementation in CAMB.
\textit{Physical Review D}, 89(10), 103530.
\href{https://doi.org/10.1103/PhysRevD.89.103530}{DOI: 10.1103/PhysRevD.89.103530}. Code: \url{https://github.com/EFTCAMB/EFTCAMB}.

\bibitem{Pitrou2018}
Pitrou, C., Coc, A., Uzan, J.\,P.\ \& Vangioni, E.\ (2018).
Precision Big Bang Nucleosynthesis with Improved Helium-4 Predictions.
\textit{Physics Reports}, 754, 1--66.
\href{https://doi.org/10.1016/j.physrep.2018.04.005}{DOI: 10.1016/j.physrep.2018.04.005}.

\bibitem{Alam2017}
Alam, S.\ et al.\ (BOSS Collaboration) (2017).
The Clustering of Galaxies in the Completed SDSS-III Baryon Oscillation Spectroscopic Survey.
\textit{Monthly Notices of the Royal Astronomical Society}, 470(3), 2617--2652.
\href{https://doi.org/10.1093/mnras/stx721}{DOI: 10.1093/mnras/stx721}.

\bibitem{Abbott2017gw}
Abbott, B.\,P.\ et al.\ (LIGO/Virgo Collaboration) (2017).
GW170817: Observation of Gravitational Waves from a Binary Neutron Star Inspiral.
\textit{Physical Review Letters}, 119(16), 161101.
\href{https://doi.org/10.1103/PhysRevLett.119.161101}{DOI: 10.1103/PhysRevLett.119.161101}.

\bibitem{Amendola2018}
Amendola, L., Kunz, M., Saltas, I.\,D.\ \& Sawicki, I.\ (2018).
Fate of Large-Scale Structure in Modified Gravity After GW170817 and GRB170817A.
\textit{Physical Review Letters}, 120(13), 131101.
\href{https://doi.org/10.1103/PhysRevLett.120.131101}{DOI: 10.1103/PhysRevLett.120.131101}.

\bibitem{Blas2011}
Blas, D., Lesgourgues, J.\ \& Tram, T.\ (2011).
The Cosmic Linear Anisotropy Solving System (CLASS). Part~II: Approximation schemes.
\textit{Journal of Cosmology and Astroparticle Physics}, 2011(07), 034.
\href{https://doi.org/10.1088/1475-7516/2011/07/034}{DOI: 10.1088/1475-7516/2011/07/034}.

\bibitem{Harris2020}
Harris, C.\,R.\ et al.\ (2020).
Array programming with NumPy.
\textit{Nature}, 585, 357--362.
\href{https://doi.org/10.1038/s41586-020-2649-2}{DOI: 10.1038/s41586-020-2649-2}.

\bibitem{Virtanen2020}
Virtanen, P.\ et al.\ (2020).
SciPy 1.0: Fundamental Algorithms for Scientific Computing in Python.
\textit{Nature Methods}, 17, 261--272.
\href{https://doi.org/10.1038/s41592-019-0686-2}{DOI: 10.1038/s41592-019-0686-2}.

\bibitem{Hunter2007}
Hunter, J.\,D.\ (2007).
Matplotlib: A 2D Graphics Environment.
\textit{Computing in Science \& Engineering}, 9(3), 90--95.
\href{https://doi.org/10.1109/MCSE.2007.55}{DOI: 10.1109/MCSE.2007.55}.

\bibitem{ForemanMackey2013}
Foreman-Mackey, D., Hogg, D.\,W., Lang, D.\ \& Goodman, J.\ (2013).
emcee: The MCMC Hammer.
\textit{Publications of the Astronomical Society of the Pacific}, 125(925), 306--312.
\href{https://doi.org/10.1086/670067}{DOI: 10.1086/670067}.

\bibitem{FamaeyMcGaugh2012}
B.~Famaey and S.~S.~McGaugh (2012).
Modified Newtonian Dynamics (MOND): Observational Phenomenology and Relativistic Extensions.
\textit{Living Reviews in Relativity}, \textbf{15}, 10.
\href{https://doi.org/10.12942/lrr-2012-10}{DOI: 10.12942/lrr-2012-10}.

\end{thebibliography}

\end{document}
